%\documentclass[dvipdfmx,a4paper,twoside]{jsarticle}
\documentclass[uplatex,dvipdfmx,a4paper,twoside]{jsarticle}
%%%%%%%%%% 文書毎に変更しない部分を入力しておく %%%%%%%%%%
% upLaTeXコンパイルで対応する。（pLaTeXで太字明朝ができないので。）
% 偶数ページ，奇数ページのヘッダ・フッタを使い分けるために"twoside"が必要。
%\documentclass[dvipdfmx,a4paper,twoside]{jsarticle}
%\documentclass[uplaxex,dvipdfmx,a4paper,twoside]{jsarticle}

\usepackage{ketpic,ketlayer} % ketpicとlayer環境を使うため(下の2つを要求する)
\usepackage{graphicx,xcolor} % png図を挿入するためgraphicxパッケージ　＆　color を 新しいxcolorにする。

\usepackage{amsmath,amssymb,multicol,fancyhdr,mathrsfs}  % AMSパッケージ，多段組パッケージ，ヘッダ制御の追加，筆記体の指定
%{\fboxsep=1pt\fbox{\gt\rule{0em}{0.8em}解説}} % 四角囲みはこれで対応できる。

\usepackage[T1]{fontenc} % iにアクセントを付けるのが上手くいかないので。

\setmargin{24}{24}{11}{39}% ketpicパッケージを使っているのだから，余白はketpicのコマンドで．フッタを考慮。
 \setlength{\columnsep}{6mm}
\renewcommand{\baselinestretch}{.82} % 行送りを0.814--0.829倍に詰めて、48行が収まるように調整。

%%\setmargin{24}{24}{30-vf}{20+vf}% フッタの高さ・位置の分vfだけ調整する。
%%%%%%%%%% ヘッダフッタを制御する。(2ページ以降に重要)
 \pagestyle{fancy} \cfoot{}                                 % ヘッダ・フッタの制御、ページ番号の非印刷
 \setcounter{page}{1}                                       % ページの開始番号を変更
 \renewcommand{\headrulewidth}{0pt}\lhead{}\chead{}\rhead{} % ヘッダの線を引かず、ヘッダを付けない。ページ番号もなし（\thepageを除去）
 \headheight=10mm % ここの値が，上のvfに影響を与える。
 \headsep=8mm % ヘッダと本文が重なるので挿入。

%%%%%%%%%%%%% 弓削紀要様式にするためのプリアンブル部分 %%%%%%%%%%%%%

%%%% ２ページ以降のヘッダ・フッタのためのコマンド %%%%
\newcommand{\secondhead}[2]{
\pagestyle{fancy}\renewcommand{\headrulewidth}{0pt}\renewcommand{\footrulewidth}{0pt}\lfoot{}\cfoot{}\rfoot{}
\chead[#1] {#2} }

%%%%文字コードをUTF-8へ%%%%

%%% 多書体利用の本体(明朝・ゴシックともに２つの太さが使えるようにする) %%%%
\usepackage[deluxe]{otf} % \bfseriesでもゴシックと明朝が使えるようにするため

%%%%%セクションタイトルを中央寄せにして、大きさをノーマルサイズに%%%%
% sectionタイトルの前後は2行あけ、中央寄せにする。section番号とsectionタイトルは続ける。
\makeatletter%% プリアンブルで定義する場合は必須
\renewcommand{\section}{
  \@startsection{section}% #1 見出し
   {1}% #2 見出しのレベル
   {\z@}% #3 横組みの場合，見出し左の空き(インデント量)
   {.5\Cvs \@plus .5\Cdp \@minus.2\Cdp}% #4 見出し上の空き
   {.5\Cvs \@plus.9\Cdp}% #5 見出し下の空き (負の値なら見出し後の空き) 
   {\centering\reset@font\normalsize\gtfamily\bfseries}% #6 見出しの属性(中央揃え＆ゴシック・太字)
}% renwecommandの終わり
\makeatother%% プリアンブルで定義する場合は必須
%%%%%%%%%%%%%%%%%%%%%%%%%%%%%%%%%%%%%%%%%%%%

%%%%subsectionタイトルを左寄せ、ノーマルサイズ、前は１行空け、後ろは通常の行送り%%%%
\makeatletter%% プリアンブルで定義する場合は必須
\renewcommand{\subsection}{
  \@startsection{subsection}% #1 見出し
   {2}% #2 見出しのレベル
   {\z@}% #3 横組みの場合，見出し左の空き(インデント量) カウンタ数字を含めた開始位置
   {.5\Cvs \@plus.9\Cdp \@minus.2\Cdp}% #4 見出し上の空き
   {.0\Cvs \@plus.0\Cdp}% #5 見出し下の空き (負の値なら見出し後の空き) 
  {\reset@font\normalsize\gtfamily\bfseries}% #6 見出しの属性（ゴチック&太字）
}% renwecommandの終わり
\makeatother%% プリアンブルで定義する場合は必須
%%%%%%%%%%%%%%%%%%%%%%%%%%%%%%%%%%%%%%%%%%%%55

%%%% section の番号に全角数字を使うことで、紀要の様式に合わせる %%%%
\makeatletter%% プリアンブルで定義する場合は必須
\def\@arabicz#1{%
  \ifcase#1 ０\or １\or ２\or ３\or ４\or ５\or ６\or ７\or ８\or ９\or １０\or
  １１\or １２\or １３\or １４\or １５\or １６\or １７\or １８\else\@ctrerr\fi} % (1～10を用意)
  \def\arabicz#1{\expandafter\@arabicz\csname c@#1\endcsname}%１２３ 全角
\makeatother%% プリアンブルで定義する場合は必須

%%%%%%%%%%%% 上で定義した数字を節番号に利用する %%%%%%%%%%%%%%%%%%
\renewcommand{\thesection}{\arabicz{section}{.\hspace{-1.2ex}}} % 節番号をピリオド付きの全角数字にする。
\renewcommand{\thesubsection}{\thesection{\hspace{2ex}\arabicz{subsection}}} % sub節番号をピリオド付きの全角数字にする。

\renewcommand{\thesubsubsection}{\thesubsection{\hspace{.2ex}.\arabicz{subsubsection}}} % subsub節番号をピリオド付きの全角数字にする。

%%%%%%%%%%%% 図のキャプション番号の様式を変更する %%%%%%%%%%%%%%%%%%
\makeatletter
\long\def\@makecaption#1#2{ %
\vskip\abovecaptionskip
\iftdir\sbox\@tempboxa{#1\hskip1zw#2} %
\else\sbox\@tempboxa{\textbf{#1}　#2} % ←変更
\fi
\ifdim \wd\@tempboxa >\hsize%
\iftdir #1\hskip1zw#2\relax\par
\else \textbf{#1:} #2\relax\par\fi % ←変更
\else
\global \@minipagefalse
\hbox to\hsize{\hfil\box\@tempboxa\hfil}
\fi
\vskip\belowcaptionskip}
\makeatother

%%%% 行間の指定 %%%%
%この後の行間の指定が、本文の式やsection 位置によって勝手に調整されないように次の2行を加えておく。
\raggedbottom  % 通常の対策
\raggedcolumns % multicol環境に対する対策

%%%%紀要は図の前後を2行ずつ空ける%%%%
% 2段組ページの途中に出力される図と本文の間の余白※未解決
%\abovecaptionskip=0pt \belowcaptionskip=24pt % 図のキャプションの上下を指定、
 \abovecaptionskip=0pt \belowcaptionskip=1.9\Cvs % 図のキャプションの上下を指定、紀要は図の前後を2行ずつ空けるのでキャプションの下だけ２行空ける

%%%%%%%%% 2段組の文章中に画像を挿入するための定義（以下の4行）%%%%%%%%%%
\makeatletter
\newenvironment{tablehere}  {\def\@captype{table}}  {}
\newenvironment{figurehere}  {\def\@captype{figure}}  {}
\makeatother

%%%%%%%%%% 行間の調整めやす %%%%%%%%%%
% subsectionタイトルは前のみ1行あける。1行の中で左寄せ(本文と同じ）。本文とは通常の改行。section番号とsectionタイトルは、1文字あける。
% 図の前後は2行あける。図と図も2行あける。図とsectioタイトル、図とsubsectionタイトルの間は間隔の大きい方とし、2行あける。

% 「Abstract」は、centuryの太字を指定されているので、\bf で代用する。

%%%%%%%%%% 短縮形のためのコマンド %%%%%%%%%%

% \単語Abstract
\newcommand{\単語Abstract}{{
{\bf\fontsize{12pt}{0pt}\selectfont \hfil\ {\hbox to 17mm{A\!\hfil b\!\hfil s\!\hfil t\!\hfil r\!\hfil a\!\hfil c\!\hfil t}}\hfil\ \ }
}}
 % 弓削紀要の様式(タイトル、番号)に合わせるためのもの
% upLaTeXでコンパイルするように変更（昨年までと異なる）。

\usepackage[bookmarks=false,colorlinks=true,urlcolor=blue]{hyperref}
\usepackage{url}

\usepackage[dvipdfmx]{pict2e}

\newcommand{\maintitlej}{{\ketcindy {の紹介}}} % ここで和文タイトルを定義して使い回す．

\begin{document}

%%%% 最初のページのヘッダ・フッタ(ここから) %%%% ヘッダ無し、フッタあり。(フッタ：所属と提出日)
\thispagestyle{empty}

\begin{center}
\vspace*{-6mm}
{\fontsize{21pt}{0pt}\selectfont \textbf{\maintitlej}} % 和文表題を\maintitlejで挿入

% {\fontsize{16pt}{0pt}\selectfont --- 平成３０年度の準備(副題16pt) ---} % 和文副題
 \end{center}
\vspace*{.3\baselineskip} % centring による行送りを考慮して、10ポイントの１行空きに調整する。

{\setlength{\baselineskip}{19pt}
\begin{center} {\fontsize{17pt}{0pt}\selectfont\bf 
--- 点や直線を動かしてみよう ---
} % 英文タイトル

%{\fontsize{14pt}{0pt}\selectfont --- Sub Title(14point, centering) ---} % 英文サブタイトル
\end{center} }
\vspace*{.3\baselineskip}


\vspace*{2\baselineskip}
%%%% ここまででタイトル・著者・Abstract など本文に入る前の準備をすます。%%%%

\begin{multicols}{2}
\vspace*{-2.5\baselineskip} % 2段組の開始がタイトルなので、右の段に会わせるための調整
\section{直線をはさむ2点}

\subsection{2点P,Q は対称}

直線をはさむ2点を見てみよう。\par
\begin{layer}{80}{0}
\putnotese{7}{7}{%%% /Users/kumac21/DataKubo_Mac21/予定とメモなど/予定とメモ（学内）/予定：学内21/学事予定21/オープンキャンパス21秋/出品(久保)/fig/対称移動1.tex 
%%% Generator=対称移動1.cdy 
{\unitlength=8mm%
\begin{picture}%
(7.45,7.45)(-2.45,-2.45)%
\special{pn 8}%
%
\special{pa  -772   772}\special{pa  -756   756}\special{pa  -693   693}\special{pa  -630   630}%
\special{pa  -567   567}\special{pa  -504   504}\special{pa  -441   441}\special{pa  -378   378}%
\special{pa  -315   315}\special{pa  -252   252}\special{pa  -189   189}\special{pa  -126   126}%
\special{pa   -63    63}\special{pa     0    -0}\special{pa    63   -63}\special{pa   126  -126}%
\special{pa   189  -189}\special{pa   252  -252}\special{pa   315  -315}\special{pa   378  -378}%
\special{pa   441  -441}\special{pa   504  -504}\special{pa   567  -567}\special{pa   630  -630}%
\special{pa   693  -693}\special{pa   756  -756}\special{pa   819  -819}\special{pa   882  -882}%
\special{pa   945  -945}\special{pa  1008 -1008}\special{pa  1071 -1071}\special{pa  1134 -1134}%
\special{pa  1197 -1197}\special{pa  1260 -1260}\special{pa  1323 -1323}\special{pa  1386 -1386}%
\special{pa  1449 -1449}\special{pa  1512 -1512}\special{pa  1575 -1575}%
\special{fp}%
\special{pa 951 -315}\special{pa 951 -316}\special{pa 951 -316}\special{pa 950 -317}%
\special{pa 950 -318}\special{pa 950 -318}\special{pa 949 -319}\special{pa 949 -320}%
\special{pa 948 -320}\special{pa 947 -320}\special{pa 947 -321}\special{pa 946 -321}%
\special{pa 945 -321}\special{pa 945 -321}\special{pa 944 -321}\special{pa 943 -321}%
\special{pa 942 -320}\special{pa 942 -320}\special{pa 941 -320}\special{pa 941 -319}%
\special{pa 940 -318}\special{pa 940 -318}\special{pa 939 -317}\special{pa 939 -316}%
\special{pa 939 -316}\special{pa 939 -315}\special{pa 939 -314}\special{pa 939 -313}%
\special{pa 939 -313}\special{pa 940 -312}\special{pa 940 -311}\special{pa 941 -311}%
\special{pa 941 -310}\special{pa 942 -310}\special{pa 942 -310}\special{pa 943 -309}%
\special{pa 944 -309}\special{pa 945 -309}\special{pa 945 -309}\special{pa 946 -309}%
\special{pa 947 -309}\special{pa 947 -310}\special{pa 948 -310}\special{pa 949 -310}%
\special{pa 949 -311}\special{pa 950 -311}\special{pa 950 -312}\special{pa 950 -313}%
\special{pa 951 -313}\special{pa 951 -314}\special{pa 951 -315}\special{pa 951 -315}%
\special{sh 1}\special{ip}%
\special{pa   951  -315}\special{pa   951  -316}\special{pa   951  -316}\special{pa   950  -317}%
\special{pa   950  -318}\special{pa   950  -318}\special{pa   949  -319}\special{pa   949  -320}%
\special{pa   948  -320}\special{pa   947  -320}\special{pa   947  -321}\special{pa   946  -321}%
\special{pa   945  -321}\special{pa   945  -321}\special{pa   944  -321}\special{pa   943  -321}%
\special{pa   942  -320}\special{pa   942  -320}\special{pa   941  -320}\special{pa   941  -319}%
\special{pa   940  -318}\special{pa   940  -318}\special{pa   939  -317}\special{pa   939  -316}%
\special{pa   939  -316}\special{pa   939  -315}\special{pa   939  -314}\special{pa   939  -313}%
\special{pa   939  -313}\special{pa   940  -312}\special{pa   940  -311}\special{pa   941  -311}%
\special{pa   941  -310}\special{pa   942  -310}\special{pa   942  -310}\special{pa   943  -309}%
\special{pa   944  -309}\special{pa   945  -309}\special{pa   945  -309}\special{pa   946  -309}%
\special{pa   947  -309}\special{pa   947  -310}\special{pa   948  -310}\special{pa   949  -310}%
\special{pa   949  -311}\special{pa   950  -311}\special{pa   950  -312}\special{pa   950  -313}%
\special{pa   951  -313}\special{pa   951  -314}\special{pa   951  -315}%
\special{fp}%
\special{pa 321 -945}\special{pa 321 -946}\special{pa 321 -946}\special{pa 321 -947}%
\special{pa 320 -948}\special{pa 320 -948}\special{pa 319 -949}\special{pa 319 -949}%
\special{pa 318 -950}\special{pa 318 -950}\special{pa 317 -951}\special{pa 316 -951}%
\special{pa 315 -951}\special{pa 315 -951}\special{pa 314 -951}\special{pa 313 -951}%
\special{pa 312 -950}\special{pa 312 -950}\special{pa 311 -949}\special{pa 311 -949}%
\special{pa 310 -948}\special{pa 310 -948}\special{pa 309 -947}\special{pa 309 -946}%
\special{pa 309 -946}\special{pa 309 -945}\special{pa 309 -944}\special{pa 309 -943}%
\special{pa 309 -943}\special{pa 310 -942}\special{pa 310 -941}\special{pa 311 -941}%
\special{pa 311 -940}\special{pa 312 -940}\special{pa 312 -939}\special{pa 313 -939}%
\special{pa 314 -939}\special{pa 315 -939}\special{pa 315 -939}\special{pa 316 -939}%
\special{pa 317 -939}\special{pa 318 -939}\special{pa 318 -940}\special{pa 319 -940}%
\special{pa 319 -941}\special{pa 320 -941}\special{pa 320 -942}\special{pa 321 -943}%
\special{pa 321 -943}\special{pa 321 -944}\special{pa 321 -945}\special{pa 321 -945}%
\special{sh 1}\special{ip}%
\special{pa   321  -945}\special{pa   321  -946}\special{pa   321  -946}\special{pa   321  -947}%
\special{pa   320  -948}\special{pa   320  -948}\special{pa   319  -949}\special{pa   319  -949}%
\special{pa   318  -950}\special{pa   318  -950}\special{pa   317  -951}\special{pa   316  -951}%
\special{pa   315  -951}\special{pa   315  -951}\special{pa   314  -951}\special{pa   313  -951}%
\special{pa   312  -950}\special{pa   312  -950}\special{pa   311  -949}\special{pa   311  -949}%
\special{pa   310  -948}\special{pa   310  -948}\special{pa   309  -947}\special{pa   309  -946}%
\special{pa   309  -946}\special{pa   309  -945}\special{pa   309  -944}\special{pa   309  -943}%
\special{pa   309  -943}\special{pa   310  -942}\special{pa   310  -941}\special{pa   311  -941}%
\special{pa   311  -940}\special{pa   312  -940}\special{pa   312  -939}\special{pa   313  -939}%
\special{pa   314  -939}\special{pa   315  -939}\special{pa   315  -939}\special{pa   316  -939}%
\special{pa   317  -939}\special{pa   318  -939}\special{pa   318  -940}\special{pa   319  -940}%
\special{pa   319  -941}\special{pa   320  -941}\special{pa   320  -942}\special{pa   321  -943}%
\special{pa   321  -943}\special{pa   321  -944}\special{pa   321  -945}%
\special{fp}%
\special{pa   945  -315}\special{pa   315  -945}%
\special{fp}%
\special{pa   776  -450}\special{pa   810  -484}%
\special{fp}%
\special{pa   765  -461}\special{pa   799  -495}%
\special{fp}%
\special{pa   484  -810}\special{pa   450  -776}%
\special{fp}%
\special{pa   495  -799}\special{pa   461  -765}%
\special{fp}%
\special{pa  -315   -20}\special{pa  -315    20}%
\special{fp}%
\settowidth{\Width}{$1$}\setlength{\Width}{-0.5\Width}%
\settoheight{\Height}{$1$}\settodepth{\Depth}{$1$}\setlength{\Height}{-\Height}%
\put(1.0000000,-0.1250000){\hspace*{\Width}\raisebox{\Height}{$1$}}%
%
\special{pa  -630   -20}\special{pa  -630    20}%
\special{fp}%
\settowidth{\Width}{$2$}\setlength{\Width}{-0.5\Width}%
\settoheight{\Height}{$2$}\settodepth{\Depth}{$2$}\setlength{\Height}{-\Height}%
\put(2.0000000,-0.1250000){\hspace*{\Width}\raisebox{\Height}{$2$}}%
%
\special{pa   945   -20}\special{pa   945    20}%
\special{fp}%
\settowidth{\Width}{$3$}\setlength{\Width}{-0.5\Width}%
\settoheight{\Height}{$3$}\settodepth{\Depth}{$3$}\setlength{\Height}{-\Height}%
\put(3.0000000,-0.1250000){\hspace*{\Width}\raisebox{\Height}{$3$}}%
%
\special{pa  1260   -20}\special{pa  1260    20}%
\special{fp}%
\settowidth{\Width}{$4$}\setlength{\Width}{-0.5\Width}%
\settoheight{\Height}{$4$}\settodepth{\Depth}{$4$}\setlength{\Height}{-\Height}%
\put(4.0000000,-0.1250000){\hspace*{\Width}\raisebox{\Height}{$4$}}%
%
\special{pa  -315   -20}\special{pa  -315    20}%
\special{fp}%
\settowidth{\Width}{$-1$}\setlength{\Width}{-0.5\Width}%
\settoheight{\Height}{$-1$}\settodepth{\Depth}{$-1$}\setlength{\Height}{-\Height}%
\put(-1.0000000,-0.1250000){\hspace*{\Width}\raisebox{\Height}{$-1$}}%
%
\special{pa  -630   -20}\special{pa  -630    20}%
\special{fp}%
\settowidth{\Width}{$-2$}\setlength{\Width}{-0.5\Width}%
\settoheight{\Height}{$-2$}\settodepth{\Depth}{$-2$}\setlength{\Height}{-\Height}%
\put(-2.0000000,-0.1250000){\hspace*{\Width}\raisebox{\Height}{$-2$}}%
%
\special{pa    20   315}\special{pa   -20   315}%
\special{fp}%
\settowidth{\Width}{$1$}\setlength{\Width}{-1\Width}%
\settoheight{\Height}{$1$}\settodepth{\Depth}{$1$}\setlength{\Height}{-0.5\Height}\setlength{\Depth}{0.5\Depth}\addtolength{\Height}{\Depth}%
\put(-0.1250000,1.0000000){\hspace*{\Width}\raisebox{\Height}{$1$}}%
%
\special{pa    20   630}\special{pa   -20   630}%
\special{fp}%
\settowidth{\Width}{$2$}\setlength{\Width}{-1\Width}%
\settoheight{\Height}{$2$}\settodepth{\Depth}{$2$}\setlength{\Height}{-0.5\Height}\setlength{\Depth}{0.5\Depth}\addtolength{\Height}{\Depth}%
\put(-0.1250000,2.0000000){\hspace*{\Width}\raisebox{\Height}{$2$}}%
%
\special{pa    20  -945}\special{pa   -20  -945}%
\special{fp}%
\settowidth{\Width}{$3$}\setlength{\Width}{-1\Width}%
\settoheight{\Height}{$3$}\settodepth{\Depth}{$3$}\setlength{\Height}{-0.5\Height}\setlength{\Depth}{0.5\Depth}\addtolength{\Height}{\Depth}%
\put(-0.1250000,3.0000000){\hspace*{\Width}\raisebox{\Height}{$3$}}%
%
\special{pa    20 -1260}\special{pa   -20 -1260}%
\special{fp}%
\settowidth{\Width}{$4$}\setlength{\Width}{-1\Width}%
\settoheight{\Height}{$4$}\settodepth{\Depth}{$4$}\setlength{\Height}{-0.5\Height}\setlength{\Depth}{0.5\Depth}\addtolength{\Height}{\Depth}%
\put(-0.1250000,4.0000000){\hspace*{\Width}\raisebox{\Height}{$4$}}%
%
\special{pa    20   315}\special{pa   -20   315}%
\special{fp}%
\settowidth{\Width}{$-1$}\setlength{\Width}{-1\Width}%
\settoheight{\Height}{$-1$}\settodepth{\Depth}{$-1$}\setlength{\Height}{-0.5\Height}\setlength{\Depth}{0.5\Depth}\addtolength{\Height}{\Depth}%
\put(-0.1250000,-1.0000000){\hspace*{\Width}\raisebox{\Height}{$-1$}}%
%
\special{pa    20   630}\special{pa   -20   630}%
\special{fp}%
\settowidth{\Width}{$-2$}\setlength{\Width}{-1\Width}%
\settoheight{\Height}{$-2$}\settodepth{\Depth}{$-2$}\setlength{\Height}{-0.5\Height}\setlength{\Depth}{0.5\Depth}\addtolength{\Height}{\Depth}%
\put(-0.1250000,-2.0000000){\hspace*{\Width}\raisebox{\Height}{$-2$}}%
%
\settowidth{\Width}{P}\setlength{\Width}{0\Width}%
\settoheight{\Height}{P}\settodepth{\Depth}{P}\setlength{\Height}{-0.5\Height}\setlength{\Depth}{0.5\Depth}\addtolength{\Height}{\Depth}%
\put(3.1875000,1.0000000){\hspace*{\Width}\raisebox{\Height}{P}}%
%
\settowidth{\Width}{Q}\setlength{\Width}{0\Width}%
\settoheight{\Height}{Q}\settodepth{\Depth}{Q}\setlength{\Height}{\Depth}%
\put(1.0625000,3.1875000){\hspace*{\Width}\raisebox{\Height}{Q}}%
%
\special{pa 1500 24}\special{pa 1575 0}\special{pa 1500 -24}\special{pa 1515 0}\special{pa 1500 24}%
\special{pa 1500 24}\special{sh 1}\special{ip}%
\special{pn 1}%
\special{pa  1500    24}\special{pa  1575    -0}\special{pa  1500   -24}\special{pa  1515    -0}%
\special{pa  1500    24}%
\special{fp}%
\special{pn 8}%
\special{pn 8}%
\special{pa  -772    -0}\special{pa  1515    -0}%
\special{fp}%
\special{pn 8}%
\special{pa 24 -1500}\special{pa 0 -1575}\special{pa -24 -1500}\special{pa 0 -1515}%
\special{pa 24 -1500}\special{pa 24 -1500}\special{sh 1}\special{ip}%
\special{pn 1}%
\special{pa    24 -1500}\special{pa     0 -1575}\special{pa   -24 -1500}\special{pa     0 -1515}%
\special{pa    24 -1500}%
\special{fp}%
\special{pn 8}%
\special{pn 8}%
\special{pa     0   772}\special{pa     0 -1515}%
\special{fp}%
\special{pn 8}%
\settowidth{\Width}{$x$}\setlength{\Width}{0\Width}%
\settoheight{\Height}{$x$}\settodepth{\Depth}{$x$}\setlength{\Height}{-0.5\Height}\setlength{\Depth}{0.5\Depth}\addtolength{\Height}{\Depth}%
\put(5.0625000,0.0000000){\hspace*{\Width}\raisebox{\Height}{$x$}}%
%
\settowidth{\Width}{$y$}\setlength{\Width}{-0.5\Width}%
\settoheight{\Height}{$y$}\settodepth{\Depth}{$y$}\setlength{\Height}{\Depth}%
\put(0.0000000,5.0625000){\hspace*{\Width}\raisebox{\Height}{$y$}}%
%
\settowidth{\Width}{O}\setlength{\Width}{0\Width}%
\settoheight{\Height}{O}\settodepth{\Depth}{O}\setlength{\Height}{-\Height}%
\put(0.0625000,-0.0625000){\hspace*{\Width}\raisebox{\Height}{O}}%
%
\end{picture}}%}
\end{layer}\par
\vspace{15\baselineskip}

下のQRコードからリンクへ飛んでみて。\\
点Pをつまんで動かすと，点Qも動くよ。\\

%% upLaTeXでは，図の挿入にBBdataが必須
\includegraphics[bb=18 18 170 170,width=50mm,clip]{対称移動1QR.png}\\

\url{https://kubo-yuge.github.io/class2019/taisyoidou1/taisyoidou1json.html}\\

\columnbreak

\subsection{別の直線でも}

直線をはさむ2点を見てみよう。\par
\begin{layer}{80}{0}
\putnotese{7}{7}{%%% /Users/kumac21/DataKubo_Mac21/予定とメモなど/予定とメモ（学内）/予定：学内21/学事予定21/オープンキャンパス21秋/出品(久保)/fig/対称移動2.tex 
%%% Generator=対称移動2.cdy 
{\unitlength=8mm%
\begin{picture}%
(7.45,7.45)(-2.45,-2.45)%
\special{pn 8}%
%
\special{pa  -386   772}\special{pa  -378   756}\special{pa  -315   630}\special{pa  -252   504}%
\special{pa  -189   378}\special{pa  -126   252}\special{pa   -63   126}\special{pa     0    -0}%
\special{pa    63  -126}\special{pa   126  -252}\special{pa   189  -378}\special{pa   252  -504}%
\special{pa   315  -630}\special{pa   378  -756}\special{pa   441  -882}\special{pa   504 -1008}%
\special{pa   567 -1134}\special{pa   630 -1260}\special{pa   693 -1386}\special{pa   756 -1512}%
\special{pa   787 -1575}%
\special{fp}%
\special{pa 951 -315}\special{pa 951 -316}\special{pa 951 -316}\special{pa 950 -317}%
\special{pa 950 -318}\special{pa 950 -318}\special{pa 949 -319}\special{pa 949 -320}%
\special{pa 948 -320}\special{pa 947 -320}\special{pa 947 -321}\special{pa 946 -321}%
\special{pa 945 -321}\special{pa 945 -321}\special{pa 944 -321}\special{pa 943 -321}%
\special{pa 942 -320}\special{pa 942 -320}\special{pa 941 -320}\special{pa 941 -319}%
\special{pa 940 -318}\special{pa 940 -318}\special{pa 939 -317}\special{pa 939 -316}%
\special{pa 939 -316}\special{pa 939 -315}\special{pa 939 -314}\special{pa 939 -313}%
\special{pa 939 -313}\special{pa 940 -312}\special{pa 940 -311}\special{pa 941 -311}%
\special{pa 941 -310}\special{pa 942 -310}\special{pa 942 -310}\special{pa 943 -309}%
\special{pa 944 -309}\special{pa 945 -309}\special{pa 945 -309}\special{pa 946 -309}%
\special{pa 947 -309}\special{pa 947 -310}\special{pa 948 -310}\special{pa 949 -310}%
\special{pa 949 -311}\special{pa 950 -311}\special{pa 950 -312}\special{pa 950 -313}%
\special{pa 951 -313}\special{pa 951 -314}\special{pa 951 -315}\special{pa 951 -315}%
\special{sh 1}\special{ip}%
\special{pa   951  -315}\special{pa   951  -316}\special{pa   951  -316}\special{pa   950  -317}%
\special{pa   950  -318}\special{pa   950  -318}\special{pa   949  -319}\special{pa   949  -320}%
\special{pa   948  -320}\special{pa   947  -320}\special{pa   947  -321}\special{pa   946  -321}%
\special{pa   945  -321}\special{pa   945  -321}\special{pa   944  -321}\special{pa   943  -321}%
\special{pa   942  -320}\special{pa   942  -320}\special{pa   941  -320}\special{pa   941  -319}%
\special{pa   940  -318}\special{pa   940  -318}\special{pa   939  -317}\special{pa   939  -316}%
\special{pa   939  -316}\special{pa   939  -315}\special{pa   939  -314}\special{pa   939  -313}%
\special{pa   939  -313}\special{pa   940  -312}\special{pa   940  -311}\special{pa   941  -311}%
\special{pa   941  -310}\special{pa   942  -310}\special{pa   942  -310}\special{pa   943  -309}%
\special{pa   944  -309}\special{pa   945  -309}\special{pa   945  -309}\special{pa   946  -309}%
\special{pa   947  -309}\special{pa   947  -310}\special{pa   948  -310}\special{pa   949  -310}%
\special{pa   949  -311}\special{pa   950  -311}\special{pa   950  -312}\special{pa   950  -313}%
\special{pa   951  -313}\special{pa   951  -314}\special{pa   951  -315}%
\special{fp}%
\special{pa -309 -945}\special{pa -309 -946}\special{pa -309 -946}\special{pa -309 -947}%
\special{pa -310 -948}\special{pa -310 -948}\special{pa -311 -949}\special{pa -311 -949}%
\special{pa -312 -950}\special{pa -312 -950}\special{pa -313 -951}\special{pa -314 -951}%
\special{pa -315 -951}\special{pa -315 -951}\special{pa -316 -951}\special{pa -317 -951}%
\special{pa -318 -950}\special{pa -318 -950}\special{pa -319 -949}\special{pa -319 -949}%
\special{pa -320 -948}\special{pa -320 -948}\special{pa -321 -947}\special{pa -321 -946}%
\special{pa -321 -946}\special{pa -321 -945}\special{pa -321 -944}\special{pa -321 -943}%
\special{pa -321 -943}\special{pa -320 -942}\special{pa -320 -941}\special{pa -319 -941}%
\special{pa -319 -940}\special{pa -318 -940}\special{pa -318 -939}\special{pa -317 -939}%
\special{pa -316 -939}\special{pa -315 -939}\special{pa -315 -939}\special{pa -314 -939}%
\special{pa -313 -939}\special{pa -312 -939}\special{pa -312 -940}\special{pa -311 -940}%
\special{pa -311 -941}\special{pa -310 -941}\special{pa -310 -942}\special{pa -309 -943}%
\special{pa -309 -943}\special{pa -309 -944}\special{pa -309 -945}\special{pa -309 -945}%
\special{sh 1}\special{ip}%
\special{pa  -309  -945}\special{pa  -309  -946}\special{pa  -309  -946}\special{pa  -309  -947}%
\special{pa  -310  -948}\special{pa  -310  -948}\special{pa  -311  -949}\special{pa  -311  -949}%
\special{pa  -312  -950}\special{pa  -312  -950}\special{pa  -313  -951}\special{pa  -314  -951}%
\special{pa  -315  -951}\special{pa  -315  -951}\special{pa  -316  -951}\special{pa  -317  -951}%
\special{pa  -318  -950}\special{pa  -318  -950}\special{pa  -319  -949}\special{pa  -319  -949}%
\special{pa  -320  -948}\special{pa  -320  -948}\special{pa  -321  -947}\special{pa  -321  -946}%
\special{pa  -321  -946}\special{pa  -321  -945}\special{pa  -321  -944}\special{pa  -321  -943}%
\special{pa  -321  -943}\special{pa  -320  -942}\special{pa  -320  -941}\special{pa  -319  -941}%
\special{pa  -319  -940}\special{pa  -318  -940}\special{pa  -318  -939}\special{pa  -317  -939}%
\special{pa  -316  -939}\special{pa  -315  -939}\special{pa  -315  -939}\special{pa  -314  -939}%
\special{pa  -313  -939}\special{pa  -312  -939}\special{pa  -312  -940}\special{pa  -311  -940}%
\special{pa  -311  -941}\special{pa  -310  -941}\special{pa  -310  -942}\special{pa  -309  -943}%
\special{pa  -309  -943}\special{pa  -309  -944}\special{pa  -309  -945}%
\special{fp}%
\special{pa   945  -315}\special{pa  -315  -945}%
\special{fp}%
\special{pa   626  -448}\special{pa   648  -490}%
\special{fp}%
\special{pa   612  -455}\special{pa   633  -497}%
\special{fp}%
\special{pa     4  -812}\special{pa   -18  -770}%
\special{fp}%
\special{pa    18  -805}\special{pa    -4  -763}%
\special{fp}%
\special{pa  -315   -20}\special{pa  -315    20}%
\special{fp}%
\settowidth{\Width}{$1$}\setlength{\Width}{-0.5\Width}%
\settoheight{\Height}{$1$}\settodepth{\Depth}{$1$}\setlength{\Height}{-\Height}%
\put(1.0000000,-0.1250000){\hspace*{\Width}\raisebox{\Height}{$1$}}%
%
\special{pa  -630   -20}\special{pa  -630    20}%
\special{fp}%
\settowidth{\Width}{$2$}\setlength{\Width}{-0.5\Width}%
\settoheight{\Height}{$2$}\settodepth{\Depth}{$2$}\setlength{\Height}{-\Height}%
\put(2.0000000,-0.1250000){\hspace*{\Width}\raisebox{\Height}{$2$}}%
%
\special{pa   945   -20}\special{pa   945    20}%
\special{fp}%
\settowidth{\Width}{$3$}\setlength{\Width}{-0.5\Width}%
\settoheight{\Height}{$3$}\settodepth{\Depth}{$3$}\setlength{\Height}{-\Height}%
\put(3.0000000,-0.1250000){\hspace*{\Width}\raisebox{\Height}{$3$}}%
%
\special{pa  1260   -20}\special{pa  1260    20}%
\special{fp}%
\settowidth{\Width}{$4$}\setlength{\Width}{-0.5\Width}%
\settoheight{\Height}{$4$}\settodepth{\Depth}{$4$}\setlength{\Height}{-\Height}%
\put(4.0000000,-0.1250000){\hspace*{\Width}\raisebox{\Height}{$4$}}%
%
\special{pa  -315   -20}\special{pa  -315    20}%
\special{fp}%
\settowidth{\Width}{$-1$}\setlength{\Width}{-0.5\Width}%
\settoheight{\Height}{$-1$}\settodepth{\Depth}{$-1$}\setlength{\Height}{-\Height}%
\put(-1.0000000,-0.1250000){\hspace*{\Width}\raisebox{\Height}{$-1$}}%
%
\special{pa  -630   -20}\special{pa  -630    20}%
\special{fp}%
\settowidth{\Width}{$-2$}\setlength{\Width}{-0.5\Width}%
\settoheight{\Height}{$-2$}\settodepth{\Depth}{$-2$}\setlength{\Height}{-\Height}%
\put(-2.0000000,-0.1250000){\hspace*{\Width}\raisebox{\Height}{$-2$}}%
%
\special{pa    20   315}\special{pa   -20   315}%
\special{fp}%
\settowidth{\Width}{$1$}\setlength{\Width}{-1\Width}%
\settoheight{\Height}{$1$}\settodepth{\Depth}{$1$}\setlength{\Height}{-0.5\Height}\setlength{\Depth}{0.5\Depth}\addtolength{\Height}{\Depth}%
\put(-0.1250000,1.0000000){\hspace*{\Width}\raisebox{\Height}{$1$}}%
%
\special{pa    20   630}\special{pa   -20   630}%
\special{fp}%
\settowidth{\Width}{$2$}\setlength{\Width}{-1\Width}%
\settoheight{\Height}{$2$}\settodepth{\Depth}{$2$}\setlength{\Height}{-0.5\Height}\setlength{\Depth}{0.5\Depth}\addtolength{\Height}{\Depth}%
\put(-0.1250000,2.0000000){\hspace*{\Width}\raisebox{\Height}{$2$}}%
%
\special{pa    20  -945}\special{pa   -20  -945}%
\special{fp}%
\settowidth{\Width}{$3$}\setlength{\Width}{-1\Width}%
\settoheight{\Height}{$3$}\settodepth{\Depth}{$3$}\setlength{\Height}{-0.5\Height}\setlength{\Depth}{0.5\Depth}\addtolength{\Height}{\Depth}%
\put(-0.1250000,3.0000000){\hspace*{\Width}\raisebox{\Height}{$3$}}%
%
\special{pa    20 -1260}\special{pa   -20 -1260}%
\special{fp}%
\settowidth{\Width}{$4$}\setlength{\Width}{-1\Width}%
\settoheight{\Height}{$4$}\settodepth{\Depth}{$4$}\setlength{\Height}{-0.5\Height}\setlength{\Depth}{0.5\Depth}\addtolength{\Height}{\Depth}%
\put(-0.1250000,4.0000000){\hspace*{\Width}\raisebox{\Height}{$4$}}%
%
\special{pa    20   315}\special{pa   -20   315}%
\special{fp}%
\settowidth{\Width}{$-1$}\setlength{\Width}{-1\Width}%
\settoheight{\Height}{$-1$}\settodepth{\Depth}{$-1$}\setlength{\Height}{-0.5\Height}\setlength{\Depth}{0.5\Depth}\addtolength{\Height}{\Depth}%
\put(-0.1250000,-1.0000000){\hspace*{\Width}\raisebox{\Height}{$-1$}}%
%
\special{pa    20   630}\special{pa   -20   630}%
\special{fp}%
\settowidth{\Width}{$-2$}\setlength{\Width}{-1\Width}%
\settoheight{\Height}{$-2$}\settodepth{\Depth}{$-2$}\setlength{\Height}{-0.5\Height}\setlength{\Depth}{0.5\Depth}\addtolength{\Height}{\Depth}%
\put(-0.1250000,-2.0000000){\hspace*{\Width}\raisebox{\Height}{$-2$}}%
%
\settowidth{\Width}{P$(3,1)$}\setlength{\Width}{0\Width}%
\settoheight{\Height}{P$(3,1)$}\settodepth{\Depth}{P$(3,1)$}\setlength{\Height}{\Depth}%
\put(3.1875000,1.0625000){\hspace*{\Width}\raisebox{\Height}{P$(3,1)$}}%
%
\settowidth{\Width}{Q$(-1,3)$}\setlength{\Width}{-1\Width}%
\settoheight{\Height}{Q$(-1,3)$}\settodepth{\Depth}{Q$(-1,3)$}\setlength{\Height}{\Depth}%
\put(-1.0625000,3.1875000){\hspace*{\Width}\raisebox{\Height}{Q$(-1,3)$}}%
%
\settowidth{\Width}{$y=2x$}\setlength{\Width}{-0.5\Width}%
\settoheight{\Height}{$y=2x$}\settodepth{\Depth}{$y=2x$}\setlength{\Height}{\Depth}%
\put(3.0000000,4.0625000){\hspace*{\Width}\raisebox{\Height}{$y=2x$}}%
%
\special{pa 1500 24}\special{pa 1575 0}\special{pa 1500 -24}\special{pa 1515 0}\special{pa 1500 24}%
\special{pa 1500 24}\special{sh 1}\special{ip}%
\special{pn 1}%
\special{pa  1500    24}\special{pa  1575    -0}\special{pa  1500   -24}\special{pa  1515    -0}%
\special{pa  1500    24}%
\special{fp}%
\special{pn 8}%
\special{pn 8}%
\special{pa  -772    -0}\special{pa  1515    -0}%
\special{fp}%
\special{pn 8}%
\special{pa 24 -1500}\special{pa 0 -1575}\special{pa -24 -1500}\special{pa 0 -1515}%
\special{pa 24 -1500}\special{pa 24 -1500}\special{sh 1}\special{ip}%
\special{pn 1}%
\special{pa    24 -1500}\special{pa     0 -1575}\special{pa   -24 -1500}\special{pa     0 -1515}%
\special{pa    24 -1500}%
\special{fp}%
\special{pn 8}%
\special{pn 8}%
\special{pa     0   772}\special{pa     0 -1515}%
\special{fp}%
\special{pn 8}%
\settowidth{\Width}{$x$}\setlength{\Width}{0\Width}%
\settoheight{\Height}{$x$}\settodepth{\Depth}{$x$}\setlength{\Height}{-0.5\Height}\setlength{\Depth}{0.5\Depth}\addtolength{\Height}{\Depth}%
\put(5.0625000,0.0000000){\hspace*{\Width}\raisebox{\Height}{$x$}}%
%
\settowidth{\Width}{$y$}\setlength{\Width}{-0.5\Width}%
\settoheight{\Height}{$y$}\settodepth{\Depth}{$y$}\setlength{\Height}{\Depth}%
\put(0.0000000,5.0625000){\hspace*{\Width}\raisebox{\Height}{$y$}}%
%
\settowidth{\Width}{O}\setlength{\Width}{0\Width}%
\settoheight{\Height}{O}\settodepth{\Depth}{O}\setlength{\Height}{-\Height}%
\put(0.0625000,-0.0625000){\hspace*{\Width}\raisebox{\Height}{O}}%
%
\end{picture}}%}
\end{layer}\par
\vspace{15\baselineskip}

下のQRコードからリンクへ飛んでみて。\\
点Pをつまんで動かすと，点Qも動くよ。\\

%% upLaTeXでは，図の挿入にBBdataが必須
\includegraphics[bb=18 18 170 170,width=50mm,clip]{対称移動2QR.png}\\

\url{https://kubo-yuge.github.io/class2019/taisyoidou2/taisyoidou2json.html}\\

\columnbreak

\subsection{点の座標を分数で}

直線をはさむ2点の座標に注目!!\par
\begin{layer}{80}{0}
\putnotese{3}{7}{%%% /Users/kumac21/DataKubo_Mac21/予定とメモなど/予定とメモ（学内）/予定：学内21/学事予定21/オープンキャンパス21秋/出品(久保)/fig/対称移動3.tex 
%%% Generator=対称移動3.cdy 
{\unitlength=6mm%
\begin{picture}%
(10,10)(-5,-5)%
\special{pn 8}%
%
{%
\color[cmyk]{1,1,0,0}%
\special{pa  -591  1181}\special{pa   591 -1181}%
\special{fp}%
}%
\special{pa  -236   -20}\special{pa  -236    20}%
\special{fp}%
\settowidth{\Width}{$1$}\setlength{\Width}{-0.5\Width}%
\settoheight{\Height}{$1$}\settodepth{\Depth}{$1$}\setlength{\Height}{-\Height}%
\put(1.0000000,-0.1666667){\hspace*{\Width}\raisebox{\Height}{$1$}}%
%
\special{pa  -472   -20}\special{pa  -472    20}%
\special{fp}%
\settowidth{\Width}{$2$}\setlength{\Width}{-0.5\Width}%
\settoheight{\Height}{$2$}\settodepth{\Depth}{$2$}\setlength{\Height}{-\Height}%
\put(2.0000000,-0.1666667){\hspace*{\Width}\raisebox{\Height}{$2$}}%
%
\special{pa  -709   -20}\special{pa  -709    20}%
\special{fp}%
\settowidth{\Width}{$3$}\setlength{\Width}{-0.5\Width}%
\settoheight{\Height}{$3$}\settodepth{\Depth}{$3$}\setlength{\Height}{-\Height}%
\put(3.0000000,-0.1666667){\hspace*{\Width}\raisebox{\Height}{$3$}}%
%
\special{pa  -945   -20}\special{pa  -945    20}%
\special{fp}%
\settowidth{\Width}{$4$}\setlength{\Width}{-0.5\Width}%
\settoheight{\Height}{$4$}\settodepth{\Depth}{$4$}\setlength{\Height}{-\Height}%
\put(4.0000000,-0.1666667){\hspace*{\Width}\raisebox{\Height}{$4$}}%
%
\special{pa  -236   -20}\special{pa  -236    20}%
\special{fp}%
\settowidth{\Width}{$-1$}\setlength{\Width}{-0.5\Width}%
\settoheight{\Height}{$-1$}\settodepth{\Depth}{$-1$}\setlength{\Height}{-\Height}%
\put(-1.0000000,-0.1666667){\hspace*{\Width}\raisebox{\Height}{$-1$}}%
%
\special{pa  -472   -20}\special{pa  -472    20}%
\special{fp}%
\settowidth{\Width}{$-2$}\setlength{\Width}{-0.5\Width}%
\settoheight{\Height}{$-2$}\settodepth{\Depth}{$-2$}\setlength{\Height}{-\Height}%
\put(-2.0000000,-0.1666667){\hspace*{\Width}\raisebox{\Height}{$-2$}}%
%
\special{pa  -709   -20}\special{pa  -709    20}%
\special{fp}%
\settowidth{\Width}{$-3$}\setlength{\Width}{-0.5\Width}%
\settoheight{\Height}{$-3$}\settodepth{\Depth}{$-3$}\setlength{\Height}{-\Height}%
\put(-3.0000000,-0.1666667){\hspace*{\Width}\raisebox{\Height}{$-3$}}%
%
\special{pa  -945   -20}\special{pa  -945    20}%
\special{fp}%
\settowidth{\Width}{$-4$}\setlength{\Width}{-0.5\Width}%
\settoheight{\Height}{$-4$}\settodepth{\Depth}{$-4$}\setlength{\Height}{-\Height}%
\put(-4.0000000,-0.1666667){\hspace*{\Width}\raisebox{\Height}{$-4$}}%
%
\special{pa    20   236}\special{pa   -20   236}%
\special{fp}%
\settowidth{\Width}{$1$}\setlength{\Width}{-1\Width}%
\settoheight{\Height}{$1$}\settodepth{\Depth}{$1$}\setlength{\Height}{-0.5\Height}\setlength{\Depth}{0.5\Depth}\addtolength{\Height}{\Depth}%
\put(-0.1666667,1.0000000){\hspace*{\Width}\raisebox{\Height}{$1$}}%
%
\special{pa    20   472}\special{pa   -20   472}%
\special{fp}%
\settowidth{\Width}{$2$}\setlength{\Width}{-1\Width}%
\settoheight{\Height}{$2$}\settodepth{\Depth}{$2$}\setlength{\Height}{-0.5\Height}\setlength{\Depth}{0.5\Depth}\addtolength{\Height}{\Depth}%
\put(-0.1666667,2.0000000){\hspace*{\Width}\raisebox{\Height}{$2$}}%
%
\special{pa    20   709}\special{pa   -20   709}%
\special{fp}%
\settowidth{\Width}{$3$}\setlength{\Width}{-1\Width}%
\settoheight{\Height}{$3$}\settodepth{\Depth}{$3$}\setlength{\Height}{-0.5\Height}\setlength{\Depth}{0.5\Depth}\addtolength{\Height}{\Depth}%
\put(-0.1666667,3.0000000){\hspace*{\Width}\raisebox{\Height}{$3$}}%
%
\special{pa    20   945}\special{pa   -20   945}%
\special{fp}%
\settowidth{\Width}{$4$}\setlength{\Width}{-1\Width}%
\settoheight{\Height}{$4$}\settodepth{\Depth}{$4$}\setlength{\Height}{-0.5\Height}\setlength{\Depth}{0.5\Depth}\addtolength{\Height}{\Depth}%
\put(-0.1666667,4.0000000){\hspace*{\Width}\raisebox{\Height}{$4$}}%
%
\special{pa    20   236}\special{pa   -20   236}%
\special{fp}%
\settowidth{\Width}{$-1$}\setlength{\Width}{-1\Width}%
\settoheight{\Height}{$-1$}\settodepth{\Depth}{$-1$}\setlength{\Height}{-0.5\Height}\setlength{\Depth}{0.5\Depth}\addtolength{\Height}{\Depth}%
\put(-0.1666667,-1.0000000){\hspace*{\Width}\raisebox{\Height}{$-1$}}%
%
\special{pa    20   472}\special{pa   -20   472}%
\special{fp}%
\settowidth{\Width}{$-2$}\setlength{\Width}{-1\Width}%
\settoheight{\Height}{$-2$}\settodepth{\Depth}{$-2$}\setlength{\Height}{-0.5\Height}\setlength{\Depth}{0.5\Depth}\addtolength{\Height}{\Depth}%
\put(-0.1666667,-2.0000000){\hspace*{\Width}\raisebox{\Height}{$-2$}}%
%
\special{pa    20   709}\special{pa   -20   709}%
\special{fp}%
\settowidth{\Width}{$-3$}\setlength{\Width}{-1\Width}%
\settoheight{\Height}{$-3$}\settodepth{\Depth}{$-3$}\setlength{\Height}{-0.5\Height}\setlength{\Depth}{0.5\Depth}\addtolength{\Height}{\Depth}%
\put(-0.1666667,-3.0000000){\hspace*{\Width}\raisebox{\Height}{$-3$}}%
%
\special{pa    20   945}\special{pa   -20   945}%
\special{fp}%
\settowidth{\Width}{$-4$}\setlength{\Width}{-1\Width}%
\settoheight{\Height}{$-4$}\settodepth{\Depth}{$-4$}\setlength{\Height}{-0.5\Height}\setlength{\Depth}{0.5\Depth}\addtolength{\Height}{\Depth}%
\put(-0.1666667,-4.0000000){\hspace*{\Width}\raisebox{\Height}{$-4$}}%
%
\special{pa -204 -841}\special{pa -204 -843}\special{pa -204 -844}\special{pa -205 -846}%
\special{pa -206 -847}\special{pa -206 -849}\special{pa -208 -850}\special{pa -209 -851}%
\special{pa -210 -852}\special{pa -212 -853}\special{pa -213 -854}\special{pa -215 -854}%
\special{pa -216 -854}\special{pa -218 -854}\special{pa -220 -854}\special{pa -221 -854}%
\special{pa -223 -853}\special{pa -225 -852}\special{pa -226 -851}\special{pa -227 -850}%
\special{pa -228 -849}\special{pa -229 -847}\special{pa -230 -846}\special{pa -230 -844}%
\special{pa -231 -843}\special{pa -231 -841}\special{pa -231 -839}\special{pa -230 -838}%
\special{pa -230 -836}\special{pa -229 -834}\special{pa -228 -833}\special{pa -227 -832}%
\special{pa -226 -831}\special{pa -225 -830}\special{pa -223 -829}\special{pa -221 -828}%
\special{pa -220 -828}\special{pa -218 -828}\special{pa -216 -828}\special{pa -215 -828}%
\special{pa -213 -828}\special{pa -212 -829}\special{pa -210 -830}\special{pa -209 -831}%
\special{pa -208 -832}\special{pa -206 -833}\special{pa -206 -834}\special{pa -205 -836}%
\special{pa -204 -838}\special{pa -204 -839}\special{pa -204 -841}\special{pa -204 -841}%
\special{sh 1}\special{ip}%
\special{pa  -204  -841}\special{pa  -204  -843}\special{pa  -204  -844}\special{pa  -205  -846}%
\special{pa  -206  -847}\special{pa  -206  -849}\special{pa  -208  -850}\special{pa  -209  -851}%
\special{pa  -210  -852}\special{pa  -212  -853}\special{pa  -213  -854}\special{pa  -215  -854}%
\special{pa  -216  -854}\special{pa  -218  -854}\special{pa  -220  -854}\special{pa  -221  -854}%
\special{pa  -223  -853}\special{pa  -225  -852}\special{pa  -226  -851}\special{pa  -227  -850}%
\special{pa  -228  -849}\special{pa  -229  -847}\special{pa  -230  -846}\special{pa  -230  -844}%
\special{pa  -231  -843}\special{pa  -231  -841}\special{pa  -231  -839}\special{pa  -230  -838}%
\special{pa  -230  -836}\special{pa  -229  -834}\special{pa  -228  -833}\special{pa  -227  -832}%
\special{pa  -226  -831}\special{pa  -225  -830}\special{pa  -223  -829}\special{pa  -221  -828}%
\special{pa  -220  -828}\special{pa  -218  -828}\special{pa  -216  -828}\special{pa  -215  -828}%
\special{pa  -213  -828}\special{pa  -212  -829}\special{pa  -210  -830}\special{pa  -209  -831}%
\special{pa  -208  -832}\special{pa  -206  -833}\special{pa  -206  -834}\special{pa  -205  -836}%
\special{pa  -204  -838}\special{pa  -204  -839}\special{pa  -204  -841}%
\special{fp}%
\special{pa   803  -331}\special{pa  -217  -841}%
\special{fp}%
{%
\color[cmyk]{1,1,0,0}%
\special{pa 306 -586}\special{pa 306 -588}\special{pa 306 -589}\special{pa 305 -591}%
\special{pa 305 -592}\special{pa 304 -594}\special{pa 303 -595}\special{pa 301 -596}%
\special{pa 300 -597}\special{pa 299 -598}\special{pa 297 -599}\special{pa 295 -599}%
\special{pa 294 -599}\special{pa 292 -599}\special{pa 290 -599}\special{pa 289 -599}%
\special{pa 287 -598}\special{pa 286 -597}\special{pa 284 -596}\special{pa 283 -595}%
\special{pa 282 -594}\special{pa 281 -592}\special{pa 280 -591}\special{pa 280 -589}%
\special{pa 280 -588}\special{pa 279 -586}\special{pa 280 -584}\special{pa 280 -582}%
\special{pa 280 -581}\special{pa 281 -579}\special{pa 282 -578}\special{pa 283 -577}%
\special{pa 284 -575}\special{pa 286 -574}\special{pa 287 -574}\special{pa 289 -573}%
\special{pa 290 -573}\special{pa 292 -572}\special{pa 294 -572}\special{pa 295 -573}%
\special{pa 297 -573}\special{pa 299 -574}\special{pa 300 -574}\special{pa 301 -575}%
\special{pa 303 -577}\special{pa 304 -578}\special{pa 305 -579}\special{pa 305 -581}%
\special{pa 306 -582}\special{pa 306 -584}\special{pa 306 -586}\special{pa 306 -586}%
\special{sh 1}\special{ip}%
}%
\special{pa   306  -586}\special{pa   306  -588}\special{pa   306  -589}\special{pa   305  -591}%
\special{pa   305  -592}\special{pa   304  -594}\special{pa   303  -595}\special{pa   301  -596}%
\special{pa   300  -597}\special{pa   299  -598}\special{pa   297  -599}\special{pa   295  -599}%
\special{pa   294  -599}\special{pa   292  -599}\special{pa   290  -599}\special{pa   289  -599}%
\special{pa   287  -598}\special{pa   286  -597}\special{pa   284  -596}\special{pa   283  -595}%
\special{pa   282  -594}\special{pa   281  -592}\special{pa   280  -591}\special{pa   280  -589}%
\special{pa   280  -588}\special{pa   279  -586}\special{pa   280  -584}\special{pa   280  -582}%
\special{pa   280  -581}\special{pa   281  -579}\special{pa   282  -578}\special{pa   283  -577}%
\special{pa   284  -575}\special{pa   286  -574}\special{pa   287  -574}\special{pa   289  -573}%
\special{pa   290  -573}\special{pa   292  -572}\special{pa   294  -572}\special{pa   295  -573}%
\special{pa   297  -573}\special{pa   299  -574}\special{pa   300  -574}\special{pa   301  -575}%
\special{pa   303  -577}\special{pa   304  -578}\special{pa   305  -579}\special{pa   305  -581}%
\special{pa   306  -582}\special{pa   306  -584}\special{pa   306  -586}%
\special{fp}%
\special{pa   545  -440}\special{pa   561  -471}%
\special{fp}%
\special{pa   535  -445}\special{pa   551  -477}%
\special{fp}%
\special{pa    40  -732}\special{pa    25  -700}%
\special{fp}%
\special{pa    51  -727}\special{pa    35  -695}%
\special{fp}%
\settowidth{\Width}{$\mathrm{P}\displaystyle(\frac{17}{5},\frac{7}{5})$}\setlength{\Width}{0\Width}%
\settoheight{\Height}{$\mathrm{P}\displaystyle(\frac{17}{5},\frac{7}{5})$}\settodepth{\Depth}{$\mathrm{P}\displaystyle(\frac{17}{5},\frac{7}{5})$}\setlength{\Height}{\Depth}%
\put(3.6500000,1.6500000){\hspace*{\Width}\raisebox{\Height}{$\mathrm{P}\displaystyle(\frac{17}{5},\frac{7}{5})$}}%
%
\settowidth{\Width}{$\mathrm{Q}\displaystyle(\frac{-23}{25},\frac{89}{25})$}\setlength{\Width}{-1\Width}%
\settoheight{\Height}{$\mathrm{Q}\displaystyle(\frac{-23}{25},\frac{89}{25})$}\settodepth{\Depth}{$\mathrm{Q}\displaystyle(\frac{-23}{25},\frac{89}{25})$}\setlength{\Height}{\Depth}%
\put(-1.1700000,3.8100000){\hspace*{\Width}\raisebox{\Height}{$\mathrm{Q}\displaystyle(\frac{-23}{25},\frac{89}{25})$}}%
%
\settowidth{\Width}{$y=2x$}\setlength{\Width}{0\Width}%
\settoheight{\Height}{$y=2x$}\settodepth{\Depth}{$y=2x$}\setlength{\Height}{-0.5\Height}\setlength{\Depth}{0.5\Depth}\addtolength{\Height}{\Depth}%
\put(3.0833333,5.0000000){\hspace*{\Width}\raisebox{\Height}{$y=2x$}}%
%
\special{pa 1106 24}\special{pa 1181 0}\special{pa 1106 -24}\special{pa 1121 0}\special{pa 1106 24}%
\special{pa 1106 24}\special{sh 1}\special{ip}%
\special{pn 1}%
\special{pa  1106    24}\special{pa  1181    -0}\special{pa  1106   -24}\special{pa  1121    -0}%
\special{pa  1106    24}%
\special{fp}%
\special{pn 8}%
\special{pn 8}%
\special{pa -1181    -0}\special{pa  1121    -0}%
\special{fp}%
\special{pn 8}%
\special{pa 24 -1106}\special{pa 0 -1181}\special{pa -24 -1106}\special{pa 0 -1121}%
\special{pa 24 -1106}\special{pa 24 -1106}\special{sh 1}\special{ip}%
\special{pn 1}%
\special{pa    24 -1106}\special{pa     0 -1181}\special{pa   -24 -1106}\special{pa     0 -1121}%
\special{pa    24 -1106}%
\special{fp}%
\special{pn 8}%
\special{pn 8}%
\special{pa     0  1181}\special{pa     0 -1121}%
\special{fp}%
\special{pn 8}%
\settowidth{\Width}{$x$}\setlength{\Width}{0\Width}%
\settoheight{\Height}{$x$}\settodepth{\Depth}{$x$}\setlength{\Height}{-0.5\Height}\setlength{\Depth}{0.5\Depth}\addtolength{\Height}{\Depth}%
\put(5.0833333,0.0000000){\hspace*{\Width}\raisebox{\Height}{$x$}}%
%
\settowidth{\Width}{$y$}\setlength{\Width}{-0.5\Width}%
\settoheight{\Height}{$y$}\settodepth{\Depth}{$y$}\setlength{\Height}{\Depth}%
\put(0.0000000,5.0833333){\hspace*{\Width}\raisebox{\Height}{$y$}}%
%
\settowidth{\Width}{O}\setlength{\Width}{0\Width}%
\settoheight{\Height}{O}\settodepth{\Depth}{O}\setlength{\Height}{-\Height}%
\put(0.0833333,-0.0833333){\hspace*{\Width}\raisebox{\Height}{O}}%
%
\end{picture}}%}
\end{layer}\par
\vspace{15\baselineskip}

下のQRコードからリンクへ飛んでみて。\\
点Pをつまんで動かすと，点Qも動くよ。\\

%% upLaTeXでは，図の挿入にBBdataが必須
\includegraphics[bb=18 18 185 185,width=50mm,clip]{対称移動3QR.png}\\

\url{https://kubo-yuge.github.io/class2019/taisyoidou2frac3/taisyoidou2_frac3json.html}\\

\columnbreak

\subsection{直線の傾きを変えてみよう}

白い枠の中の$m=$のあとに，整数や分数を入れると，直線の傾きが変わるよ。負の数でもいけるよ。\par
$1,\ -1,\ 2/3,\ -3/2$などでお願いします。\par
\begin{layer}{80}{0}
\putnotese{3}{7}{%%% /Users/kumac21/DataKubo_Mac21/予定とメモなど/予定とメモ（学内）/予定：学内21/学事予定21/オープンキャンパス21秋/出品(久保)/fig/対称移動4.tex 
%%% Generator=対称移動4.cdy 
{\unitlength=6mm%
\begin{picture}%
(10,10)(-5,-5)%
\special{pn 8}%
%
\settowidth{\Width}{$m=2$}\setlength{\Width}{-1\Width}%
\settoheight{\Height}{$m=2$}\settodepth{\Depth}{$m=2$}\setlength{\Height}{\Depth}%
\put(4.9166667,4.0833333){\hspace*{\Width}\raisebox{\Height}{$m=2$}}%
%
\special{pa  -591  1181}\special{pa   591 -1181}%
\special{fp}%
\special{pa  -945   -20}\special{pa  -945    20}%
\special{fp}%
\settowidth{\Width}{$-4$}\setlength{\Width}{-0.5\Width}%
\settoheight{\Height}{$-4$}\settodepth{\Depth}{$-4$}\setlength{\Height}{-\Height}%
\put(-4.0000000,-0.1666667){\hspace*{\Width}\raisebox{\Height}{$-4$}}%
%
\special{pa  -709   -20}\special{pa  -709    20}%
\special{fp}%
\settowidth{\Width}{$-3$}\setlength{\Width}{-0.5\Width}%
\settoheight{\Height}{$-3$}\settodepth{\Depth}{$-3$}\setlength{\Height}{-\Height}%
\put(-3.0000000,-0.1666667){\hspace*{\Width}\raisebox{\Height}{$-3$}}%
%
\special{pa  -472   -20}\special{pa  -472    20}%
\special{fp}%
\settowidth{\Width}{$-2$}\setlength{\Width}{-0.5\Width}%
\settoheight{\Height}{$-2$}\settodepth{\Depth}{$-2$}\setlength{\Height}{-\Height}%
\put(-2.0000000,-0.1666667){\hspace*{\Width}\raisebox{\Height}{$-2$}}%
%
\special{pa  -236   -20}\special{pa  -236    20}%
\special{fp}%
\settowidth{\Width}{$-1$}\setlength{\Width}{-0.5\Width}%
\settoheight{\Height}{$-1$}\settodepth{\Depth}{$-1$}\setlength{\Height}{-\Height}%
\put(-1.0000000,-0.1666667){\hspace*{\Width}\raisebox{\Height}{$-1$}}%
%
\special{pa   236   -20}\special{pa   236    20}%
\special{fp}%
\settowidth{\Width}{$1$}\setlength{\Width}{-0.5\Width}%
\settoheight{\Height}{$1$}\settodepth{\Depth}{$1$}\setlength{\Height}{-\Height}%
\put(1.0000000,-0.1666667){\hspace*{\Width}\raisebox{\Height}{$1$}}%
%
\special{pa   472   -20}\special{pa   472    20}%
\special{fp}%
\settowidth{\Width}{$2$}\setlength{\Width}{-0.5\Width}%
\settoheight{\Height}{$2$}\settodepth{\Depth}{$2$}\setlength{\Height}{-\Height}%
\put(2.0000000,-0.1666667){\hspace*{\Width}\raisebox{\Height}{$2$}}%
%
\special{pa   709   -20}\special{pa   709    20}%
\special{fp}%
\settowidth{\Width}{$3$}\setlength{\Width}{-0.5\Width}%
\settoheight{\Height}{$3$}\settodepth{\Depth}{$3$}\setlength{\Height}{-\Height}%
\put(3.0000000,-0.1666667){\hspace*{\Width}\raisebox{\Height}{$3$}}%
%
\special{pa   945   -20}\special{pa   945    20}%
\special{fp}%
\settowidth{\Width}{$4$}\setlength{\Width}{-0.5\Width}%
\settoheight{\Height}{$4$}\settodepth{\Depth}{$4$}\setlength{\Height}{-\Height}%
\put(4.0000000,-0.1666667){\hspace*{\Width}\raisebox{\Height}{$4$}}%
%
\special{pa    20   945}\special{pa   -20   945}%
\special{fp}%
\settowidth{\Width}{$-4$}\setlength{\Width}{-1\Width}%
\settoheight{\Height}{$-4$}\settodepth{\Depth}{$-4$}\setlength{\Height}{-0.5\Height}\setlength{\Depth}{0.5\Depth}\addtolength{\Height}{\Depth}%
\put(-0.1666667,-4.0000000){\hspace*{\Width}\raisebox{\Height}{$-4$}}%
%
\special{pa    20   709}\special{pa   -20   709}%
\special{fp}%
\settowidth{\Width}{$-3$}\setlength{\Width}{-1\Width}%
\settoheight{\Height}{$-3$}\settodepth{\Depth}{$-3$}\setlength{\Height}{-0.5\Height}\setlength{\Depth}{0.5\Depth}\addtolength{\Height}{\Depth}%
\put(-0.1666667,-3.0000000){\hspace*{\Width}\raisebox{\Height}{$-3$}}%
%
\special{pa    20   472}\special{pa   -20   472}%
\special{fp}%
\settowidth{\Width}{$-2$}\setlength{\Width}{-1\Width}%
\settoheight{\Height}{$-2$}\settodepth{\Depth}{$-2$}\setlength{\Height}{-0.5\Height}\setlength{\Depth}{0.5\Depth}\addtolength{\Height}{\Depth}%
\put(-0.1666667,-2.0000000){\hspace*{\Width}\raisebox{\Height}{$-2$}}%
%
\special{pa    20   236}\special{pa   -20   236}%
\special{fp}%
\settowidth{\Width}{$-1$}\setlength{\Width}{-1\Width}%
\settoheight{\Height}{$-1$}\settodepth{\Depth}{$-1$}\setlength{\Height}{-0.5\Height}\setlength{\Depth}{0.5\Depth}\addtolength{\Height}{\Depth}%
\put(-0.1666667,-1.0000000){\hspace*{\Width}\raisebox{\Height}{$-1$}}%
%
\special{pa    20  -236}\special{pa   -20  -236}%
\special{fp}%
\settowidth{\Width}{$1$}\setlength{\Width}{-1\Width}%
\settoheight{\Height}{$1$}\settodepth{\Depth}{$1$}\setlength{\Height}{-0.5\Height}\setlength{\Depth}{0.5\Depth}\addtolength{\Height}{\Depth}%
\put(-0.1666667,1.0000000){\hspace*{\Width}\raisebox{\Height}{$1$}}%
%
\special{pa    20  -472}\special{pa   -20  -472}%
\special{fp}%
\settowidth{\Width}{$2$}\setlength{\Width}{-1\Width}%
\settoheight{\Height}{$2$}\settodepth{\Depth}{$2$}\setlength{\Height}{-0.5\Height}\setlength{\Depth}{0.5\Depth}\addtolength{\Height}{\Depth}%
\put(-0.1666667,2.0000000){\hspace*{\Width}\raisebox{\Height}{$2$}}%
%
\special{pa    20  -709}\special{pa   -20  -709}%
\special{fp}%
\settowidth{\Width}{$3$}\setlength{\Width}{-1\Width}%
\settoheight{\Height}{$3$}\settodepth{\Depth}{$3$}\setlength{\Height}{-0.5\Height}\setlength{\Depth}{0.5\Depth}\addtolength{\Height}{\Depth}%
\put(-0.1666667,3.0000000){\hspace*{\Width}\raisebox{\Height}{$3$}}%
%
\special{pa    20  -945}\special{pa   -20  -945}%
\special{fp}%
\settowidth{\Width}{$4$}\setlength{\Width}{-1\Width}%
\settoheight{\Height}{$4$}\settodepth{\Depth}{$4$}\setlength{\Height}{-0.5\Height}\setlength{\Depth}{0.5\Depth}\addtolength{\Height}{\Depth}%
\put(-0.1666667,4.0000000){\hspace*{\Width}\raisebox{\Height}{$4$}}%
%
\special{pa 486 -236}\special{pa 486 -238}\special{pa 485 -240}\special{pa 485 -241}%
\special{pa 484 -243}\special{pa 483 -244}\special{pa 482 -245}\special{pa 481 -247}%
\special{pa 480 -248}\special{pa 478 -248}\special{pa 477 -249}\special{pa 475 -249}%
\special{pa 473 -250}\special{pa 472 -250}\special{pa 470 -249}\special{pa 468 -249}%
\special{pa 467 -248}\special{pa 465 -248}\special{pa 464 -247}\special{pa 463 -245}%
\special{pa 462 -244}\special{pa 461 -243}\special{pa 460 -241}\special{pa 459 -240}%
\special{pa 459 -238}\special{pa 459 -236}\special{pa 459 -235}\special{pa 459 -233}%
\special{pa 460 -231}\special{pa 461 -230}\special{pa 462 -228}\special{pa 463 -227}%
\special{pa 464 -226}\special{pa 465 -225}\special{pa 467 -224}\special{pa 468 -223}%
\special{pa 470 -223}\special{pa 472 -223}\special{pa 473 -223}\special{pa 475 -223}%
\special{pa 477 -223}\special{pa 478 -224}\special{pa 480 -225}\special{pa 481 -226}%
\special{pa 482 -227}\special{pa 483 -228}\special{pa 484 -230}\special{pa 485 -231}%
\special{pa 485 -233}\special{pa 486 -235}\special{pa 486 -236}\special{pa 486 -236}%
\special{sh 1}\special{ip}%
\special{pa   486  -236}\special{pa   486  -238}\special{pa   485  -240}\special{pa   485  -241}%
\special{pa   484  -243}\special{pa   483  -244}\special{pa   482  -245}\special{pa   481  -247}%
\special{pa   480  -248}\special{pa   478  -248}\special{pa   477  -249}\special{pa   475  -249}%
\special{pa   473  -250}\special{pa   472  -250}\special{pa   470  -249}\special{pa   468  -249}%
\special{pa   467  -248}\special{pa   465  -248}\special{pa   464  -247}\special{pa   463  -245}%
\special{pa   462  -244}\special{pa   461  -243}\special{pa   460  -241}\special{pa   459  -240}%
\special{pa   459  -238}\special{pa   459  -236}\special{pa   459  -235}\special{pa   459  -233}%
\special{pa   460  -231}\special{pa   461  -230}\special{pa   462  -228}\special{pa   463  -227}%
\special{pa   464  -226}\special{pa   465  -225}\special{pa   467  -224}\special{pa   468  -223}%
\special{pa   470  -223}\special{pa   472  -223}\special{pa   473  -223}\special{pa   475  -223}%
\special{pa   477  -223}\special{pa   478  -224}\special{pa   480  -225}\special{pa   481  -226}%
\special{pa   482  -227}\special{pa   483  -228}\special{pa   484  -230}\special{pa   485  -231}%
\special{pa   485  -233}\special{pa   486  -235}\special{pa   486  -236}%
\special{fp}%
\special{pa -81 -520}\special{pa -81 -521}\special{pa -81 -523}\special{pa -82 -525}%
\special{pa -83 -526}\special{pa -84 -528}\special{pa -85 -529}\special{pa -86 -530}%
\special{pa -87 -531}\special{pa -89 -532}\special{pa -90 -532}\special{pa -92 -533}%
\special{pa -94 -533}\special{pa -95 -533}\special{pa -97 -533}\special{pa -99 -532}%
\special{pa -100 -532}\special{pa -102 -531}\special{pa -103 -530}\special{pa -104 -529}%
\special{pa -105 -528}\special{pa -106 -526}\special{pa -107 -525}\special{pa -108 -523}%
\special{pa -108 -521}\special{pa -108 -520}\special{pa -108 -518}\special{pa -108 -516}%
\special{pa -107 -515}\special{pa -106 -513}\special{pa -105 -512}\special{pa -104 -510}%
\special{pa -103 -509}\special{pa -102 -508}\special{pa -100 -508}\special{pa -99 -507}%
\special{pa -97 -506}\special{pa -95 -506}\special{pa -94 -506}\special{pa -92 -506}%
\special{pa -90 -507}\special{pa -89 -508}\special{pa -87 -508}\special{pa -86 -509}%
\special{pa -85 -510}\special{pa -84 -512}\special{pa -83 -513}\special{pa -82 -515}%
\special{pa -81 -516}\special{pa -81 -518}\special{pa -81 -520}\special{pa -81 -520}%
\special{sh 1}\special{ip}%
\special{pa   -81  -520}\special{pa   -81  -521}\special{pa   -81  -523}\special{pa   -82  -525}%
\special{pa   -83  -526}\special{pa   -84  -528}\special{pa   -85  -529}\special{pa   -86  -530}%
\special{pa   -87  -531}\special{pa   -89  -532}\special{pa   -90  -532}\special{pa   -92  -533}%
\special{pa   -94  -533}\special{pa   -95  -533}\special{pa   -97  -533}\special{pa   -99  -532}%
\special{pa  -100  -532}\special{pa  -102  -531}\special{pa  -103  -530}\special{pa  -104  -529}%
\special{pa  -105  -528}\special{pa  -106  -526}\special{pa  -107  -525}\special{pa  -108  -523}%
\special{pa  -108  -521}\special{pa  -108  -520}\special{pa  -108  -518}\special{pa  -108  -516}%
\special{pa  -107  -515}\special{pa  -106  -513}\special{pa  -105  -512}\special{pa  -104  -510}%
\special{pa  -103  -509}\special{pa  -102  -508}\special{pa  -100  -508}\special{pa   -99  -507}%
\special{pa   -97  -506}\special{pa   -95  -506}\special{pa   -94  -506}\special{pa   -92  -506}%
\special{pa   -90  -507}\special{pa   -89  -508}\special{pa   -87  -508}\special{pa   -86  -509}%
\special{pa   -85  -510}\special{pa   -84  -512}\special{pa   -83  -513}\special{pa   -82  -515}%
\special{pa   -81  -516}\special{pa   -81  -518}\special{pa   -81  -520}%
\special{fp}%
\special{pa   472  -236}\special{pa   -94  -520}%
\special{fp}%
{%
\color[cmyk]{1,1,0,0}%
\special{pa 202 -378}\special{pa 202 -380}\special{pa 202 -381}\special{pa 201 -383}%
\special{pa 201 -384}\special{pa 200 -386}\special{pa 199 -387}\special{pa 198 -388}%
\special{pa 196 -389}\special{pa 195 -390}\special{pa 193 -391}\special{pa 191 -391}%
\special{pa 190 -391}\special{pa 188 -391}\special{pa 186 -391}\special{pa 185 -391}%
\special{pa 183 -390}\special{pa 182 -389}\special{pa 180 -388}\special{pa 179 -387}%
\special{pa 178 -386}\special{pa 177 -384}\special{pa 176 -383}\special{pa 176 -381}%
\special{pa 176 -380}\special{pa 176 -378}\special{pa 176 -376}\special{pa 176 -375}%
\special{pa 176 -373}\special{pa 177 -371}\special{pa 178 -370}\special{pa 179 -369}%
\special{pa 180 -368}\special{pa 182 -367}\special{pa 183 -366}\special{pa 185 -365}%
\special{pa 186 -365}\special{pa 188 -365}\special{pa 190 -365}\special{pa 191 -365}%
\special{pa 193 -365}\special{pa 195 -366}\special{pa 196 -367}\special{pa 198 -368}%
\special{pa 199 -369}\special{pa 200 -370}\special{pa 201 -371}\special{pa 201 -373}%
\special{pa 202 -375}\special{pa 202 -376}\special{pa 202 -378}\special{pa 202 -378}%
\special{sh 1}\special{ip}%
}%
\special{pa   202  -378}\special{pa   202  -380}\special{pa   202  -381}\special{pa   201  -383}%
\special{pa   201  -384}\special{pa   200  -386}\special{pa   199  -387}\special{pa   198  -388}%
\special{pa   196  -389}\special{pa   195  -390}\special{pa   193  -391}\special{pa   191  -391}%
\special{pa   190  -391}\special{pa   188  -391}\special{pa   186  -391}\special{pa   185  -391}%
\special{pa   183  -390}\special{pa   182  -389}\special{pa   180  -388}\special{pa   179  -387}%
\special{pa   178  -386}\special{pa   177  -384}\special{pa   176  -383}\special{pa   176  -381}%
\special{pa   176  -380}\special{pa   176  -378}\special{pa   176  -376}\special{pa   176  -375}%
\special{pa   176  -373}\special{pa   177  -371}\special{pa   178  -370}\special{pa   179  -369}%
\special{pa   180  -368}\special{pa   182  -367}\special{pa   183  -366}\special{pa   185  -365}%
\special{pa   186  -365}\special{pa   188  -365}\special{pa   190  -365}\special{pa   191  -365}%
\special{pa   193  -365}\special{pa   195  -366}\special{pa   196  -367}\special{pa   198  -368}%
\special{pa   199  -369}\special{pa   200  -370}\special{pa   201  -371}\special{pa   201  -373}%
\special{pa   202  -375}\special{pa   202  -376}\special{pa   202  -378}%
\special{fp}%
\special{pa   328  -289}\special{pa   344  -320}%
\special{fp}%
\special{pa   318  -294}\special{pa   333  -326}%
\special{fp}%
\special{pa    50  -467}\special{pa    34  -436}%
\special{fp}%
\special{pa    60  -462}\special{pa    45  -430}%
\special{fp}%
\settowidth{\Width}{$\mathrm{P}\displaystyle(2,1)$}\setlength{\Width}{0\Width}%
\settoheight{\Height}{$\mathrm{P}\displaystyle(2,1)$}\settodepth{\Depth}{$\mathrm{P}\displaystyle(2,1)$}\setlength{\Height}{\Depth}%
\put(2.2500000,1.2500000){\hspace*{\Width}\raisebox{\Height}{$\mathrm{P}\displaystyle(2,1)$}}%
%
\settowidth{\Width}{$\mathrm{Q}\displaystyle(\frac{-2}{5},\frac{11}{5})$}\setlength{\Width}{-1\Width}%
\settoheight{\Height}{$\mathrm{Q}\displaystyle(\frac{-2}{5},\frac{11}{5})$}\settodepth{\Depth}{$\mathrm{Q}\displaystyle(\frac{-2}{5},\frac{11}{5})$}\setlength{\Height}{\Depth}%
\put(-0.6500000,2.4500000){\hspace*{\Width}\raisebox{\Height}{$\mathrm{Q}\displaystyle(\frac{-2}{5},\frac{11}{5})$}}%
%
\settowidth{\Width}{$y=mx$}\setlength{\Width}{-1\Width}%
\settoheight{\Height}{$y=mx$}\settodepth{\Depth}{$y=mx$}\setlength{\Height}{\Depth}%
\put(4.9166667,5.0833333){\hspace*{\Width}\raisebox{\Height}{$y=mx$}}%
%
\special{pa -1181    -0}\special{pa  1181    -0}%
\special{fp}%
\special{pa     0  1181}\special{pa     0 -1181}%
\special{fp}%
\settowidth{\Width}{$x$}\setlength{\Width}{0\Width}%
\settoheight{\Height}{$x$}\settodepth{\Depth}{$x$}\setlength{\Height}{-0.5\Height}\setlength{\Depth}{0.5\Depth}\addtolength{\Height}{\Depth}%
\put(5.0833333,0.0000000){\hspace*{\Width}\raisebox{\Height}{$x$}}%
%
\settowidth{\Width}{$y$}\setlength{\Width}{-0.5\Width}%
\settoheight{\Height}{$y$}\settodepth{\Depth}{$y$}\setlength{\Height}{\Depth}%
\put(0.0000000,5.0833333){\hspace*{\Width}\raisebox{\Height}{$y$}}%
%
\settowidth{\Width}{O}\setlength{\Width}{-1\Width}%
\settoheight{\Height}{O}\settodepth{\Depth}{O}\setlength{\Height}{-\Height}%
\put(-0.0833333,-0.0833333){\hspace*{\Width}\raisebox{\Height}{O}}%
%
\end{picture}}%}
\end{layer}\par
\vspace{15\baselineskip}

下のQRコードからリンクへ飛んでみて。\\
点Pをつまんで動かすと，点Qも動くよ。\\

%% upLaTeXでは，図の挿入にBBdataが必須
\includegraphics[bb=18 18 170 170,width=50mm,clip]{対称移動4QR.png}\\

\url{https://kubo-yuge.github.io/class2019/reflection/reflectionjson.html}\\

\columnbreak

\section{どきどきハート}

スライダの点を動かすと，ハートがドキドキします。\par
\begin{layer}{80}{0}
\putnotese{3}{7}{%%% /Users/kumac21/DataKubo_Mac21/予定とメモなど/予定とメモ（学内）/予定：学内21/学事予定21/オープンキャンパス21秋/出品(久保)/fig/Heart.tex 
%%% Generator=Heart.cdy 
{\unitlength=3mm%
\begin{picture}%
(18,13)(-9,-7)%
\special{pn 8}%
%
\special{pa  -945   591}\special{pa   945   591}%
\special{fp}%
\special{pa -383 591}\special{pa -384 589}\special{pa -384 587}\special{pa -384 586}%
\special{pa -385 584}\special{pa -386 583}\special{pa -387 581}\special{pa -388 580}%
\special{pa -390 579}\special{pa -391 578}\special{pa -393 578}\special{pa -394 577}%
\special{pa -396 577}\special{pa -398 577}\special{pa -399 577}\special{pa -401 578}%
\special{pa -403 578}\special{pa -404 579}\special{pa -405 580}\special{pa -407 581}%
\special{pa -408 583}\special{pa -409 584}\special{pa -409 586}\special{pa -410 587}%
\special{pa -410 589}\special{pa -410 591}\special{pa -410 592}\special{pa -410 594}%
\special{pa -409 596}\special{pa -409 597}\special{pa -408 598}\special{pa -407 600}%
\special{pa -405 601}\special{pa -404 602}\special{pa -403 603}\special{pa -401 603}%
\special{pa -399 604}\special{pa -398 604}\special{pa -396 604}\special{pa -394 604}%
\special{pa -393 603}\special{pa -391 603}\special{pa -390 602}\special{pa -388 601}%
\special{pa -387 600}\special{pa -386 598}\special{pa -385 597}\special{pa -384 596}%
\special{pa -384 594}\special{pa -384 592}\special{pa -383 591}\special{pa -383 591}%
\special{sh 1}\special{ip}%
\special{pa  -383   591}\special{pa  -384   589}\special{pa  -384   587}\special{pa  -384   586}%
\special{pa  -385   584}\special{pa  -386   583}\special{pa  -387   581}\special{pa  -388   580}%
\special{pa  -390   579}\special{pa  -391   578}\special{pa  -393   578}\special{pa  -394   577}%
\special{pa  -396   577}\special{pa  -398   577}\special{pa  -399   577}\special{pa  -401   578}%
\special{pa  -403   578}\special{pa  -404   579}\special{pa  -405   580}\special{pa  -407   581}%
\special{pa  -408   583}\special{pa  -409   584}\special{pa  -409   586}\special{pa  -410   587}%
\special{pa  -410   589}\special{pa  -410   591}\special{pa  -410   592}\special{pa  -410   594}%
\special{pa  -409   596}\special{pa  -409   597}\special{pa  -408   598}\special{pa  -407   600}%
\special{pa  -405   601}\special{pa  -404   602}\special{pa  -403   603}\special{pa  -401   603}%
\special{pa  -399   604}\special{pa  -398   604}\special{pa  -396   604}\special{pa  -394   604}%
\special{pa  -393   603}\special{pa  -391   603}\special{pa  -390   602}\special{pa  -388   601}%
\special{pa  -387   600}\special{pa  -386   598}\special{pa  -385   597}\special{pa  -384   596}%
\special{pa  -384   594}\special{pa  -384   592}\special{pa  -383   591}%
\special{fp}%
\special{pn 24}%
\special{pa   472    -0}\special{pa   484   -15}\special{pa   495   -31}\special{pa   506   -48}%
\special{pa   517   -65}\special{pa   527   -84}\special{pa   538  -103}\special{pa   548  -122}%
\special{pa   558  -143}\special{pa   567  -165}\special{pa   576  -187}\special{pa   584  -210}%
\special{pa   591  -234}\special{pa   598  -259}\special{pa   603  -284}\special{pa   607  -309}%
\special{pa   610  -335}\special{pa   612  -362}\special{pa   611  -388}\special{pa   609  -414}%
\special{pa   605  -440}\special{pa   599  -465}\special{pa   592  -489}\special{pa   582  -513}%
\special{pa   570  -535}\special{pa   556  -556}\special{pa   540  -575}\special{pa   523  -593}%
\special{pa   503  -609}\special{pa   483  -623}\special{pa   461  -635}\special{pa   438  -645}%
\special{pa   414  -653}\special{pa   390  -659}\special{pa   365  -664}\special{pa   340  -667}%
\special{pa   315  -669}\special{pa   289  -669}\special{pa   264  -668}\special{pa   239  -665}%
\special{pa   215  -662}\special{pa   191  -658}\special{pa   168  -652}\special{pa   145  -646}%
\special{pa   122  -640}\special{pa   100  -633}\special{pa    79  -625}\special{pa    58  -617}%
\special{pa    38  -609}\special{pa    19  -600}\special{pa     0  -591}\special{pa   -19  -600}%
\special{pa   -38  -609}\special{pa   -58  -617}\special{pa   -79  -625}\special{pa  -100  -633}%
\special{pa  -122  -640}\special{pa  -145  -646}\special{pa  -168  -652}\special{pa  -191  -658}%
\special{pa  -215  -662}\special{pa  -239  -665}\special{pa  -264  -668}\special{pa  -289  -669}%
\special{pa  -315  -669}\special{pa  -340  -667}\special{pa  -365  -664}\special{pa  -390  -659}%
\special{pa  -414  -653}\special{pa  -438  -645}\special{pa  -461  -635}\special{pa  -483  -623}%
\special{pa  -503  -609}\special{pa  -523  -593}\special{pa  -540  -575}\special{pa  -556  -556}%
\special{pa  -570  -535}\special{pa  -582  -513}\special{pa  -592  -489}\special{pa  -599  -465}%
\special{pa  -605  -440}\special{pa  -609  -414}\special{pa  -611  -388}\special{pa  -612  -362}%
\special{pa  -610  -335}\special{pa  -607  -309}\special{pa  -603  -284}\special{pa  -598  -259}%
\special{pa  -591  -234}\special{pa  -584  -210}\special{pa  -576  -187}\special{pa  -567  -165}%
\special{pa  -558  -143}\special{pa  -548  -122}\special{pa  -538  -103}\special{pa  -527   -84}%
\special{pa  -517   -65}\special{pa  -506   -48}\special{pa  -495   -31}\special{pa  -484   -15}%
\special{pa  -472    -0}\special{pa  -461    14}\special{pa  -450    28}\special{pa  -439    42}%
\special{pa  -428    54}\special{pa  -417    66}\special{pa  -407    78}\special{pa  -396    88}%
\special{pa  -385    99}\special{pa  -375   109}\special{pa  -364   118}\special{pa  -354   128}%
\special{pa  -344   136}\special{pa  -334   145}\special{pa  -325   153}\special{pa  -315   160}%
\special{pa  -305   168}\special{pa  -296   175}\special{pa  -287   182}\special{pa  -277   189}%
\special{pa  -268   195}\special{pa  -259   201}\special{pa  -251   207}\special{pa  -242   213}%
\special{pa  -233   219}\special{pa  -225   225}\special{pa  -216   230}\special{pa  -208   235}%
\special{pa  -199   241}\special{pa  -191   246}\special{pa  -182   251}\special{pa  -174   256}%
\special{pa  -166   261}\special{pa  -157   266}\special{pa  -149   271}\special{pa  -141   276}%
\special{pa  -132   281}\special{pa  -124   286}\special{pa  -115   291}\special{pa  -107   296}%
\special{pa   -98   301}\special{pa   -89   306}\special{pa   -80   311}\special{pa   -71   317}%
\special{pa   -61   322}\special{pa   -52   327}\special{pa   -42   332}\special{pa   -32   338}%
\special{pa   -22   343}\special{pa   -11   349}\special{pa    -0   354}\special{pa    11   349}%
\special{pa    22   343}\special{pa    32   338}\special{pa    42   332}\special{pa    52   327}%
\special{pa    61   322}\special{pa    71   317}\special{pa    80   311}\special{pa    89   306}%
\special{pa    98   301}\special{pa   107   296}\special{pa   115   291}\special{pa   124   286}%
\special{pa   132   281}\special{pa   141   276}\special{pa   149   271}\special{pa   157   266}%
\special{pa   166   261}\special{pa   174   256}\special{pa   182   251}\special{pa   191   246}%
\special{pa   199   241}\special{pa   208   235}\special{pa   216   230}\special{pa   225   225}%
\special{pa   233   219}\special{pa   242   213}\special{pa   251   207}\special{pa   259   201}%
\special{pa   268   195}\special{pa   277   189}\special{pa   287   182}\special{pa   296   175}%
\special{pa   305   168}\special{pa   315   160}\special{pa   325   153}\special{pa   334   145}%
\special{pa   344   136}\special{pa   354   128}\special{pa   364   118}\special{pa   375   109}%
\special{pa   385    99}\special{pa   396    88}\special{pa   407    78}\special{pa   417    66}%
\special{pa   428    54}\special{pa   439    42}\special{pa   450    28}\special{pa   461    14}%
\special{pa   472     0}%
\special{fp}%
\special{pn 8}%
\settowidth{\Width}{$-2\pi$}\setlength{\Width}{-0.5\Width}%
\settoheight{\Height}{$-2\pi$}\settodepth{\Depth}{$-2\pi$}\setlength{\Height}{-\Height}%
\put(-8.0000000,-5.5000000){\hspace*{\Width}\raisebox{\Height}{$-2\pi$}}%
%
\settowidth{\Width}{$-\pi$}\setlength{\Width}{-0.5\Width}%
\settoheight{\Height}{$-\pi$}\settodepth{\Depth}{$-\pi$}\setlength{\Height}{-\Height}%
\put(-4.0000000,-5.5000000){\hspace*{\Width}\raisebox{\Height}{$-\pi$}}%
%
\settowidth{\Width}{$2\pi$}\setlength{\Width}{-0.5\Width}%
\settoheight{\Height}{$2\pi$}\settodepth{\Depth}{$2\pi$}\setlength{\Height}{-\Height}%
\put(8.0000000,-5.5000000){\hspace*{\Width}\raisebox{\Height}{$2\pi$}}%
%
\settowidth{\Width}{$\pi$}\setlength{\Width}{-0.5\Width}%
\settoheight{\Height}{$\pi$}\settodepth{\Depth}{$\pi$}\setlength{\Height}{-\Height}%
\put(4.0000000,-5.5000000){\hspace*{\Width}\raisebox{\Height}{$\pi$}}%
%
\settowidth{\Width}{$0$}\setlength{\Width}{-0.5\Width}%
\settoheight{\Height}{$0$}\settodepth{\Depth}{$0$}\setlength{\Height}{-\Height}%
\put(0.0000000,-5.5000000){\hspace*{\Width}\raisebox{\Height}{$0$}}%
%
\end{picture}}%}
\end{layer}\par
\vspace{11\baselineskip}

下のQRコードからリンクへ飛んでみて。\\

%% upLaTeXでは，図の挿入にBBdataが必須
\includegraphics[bb=10 10 220 220,width=50mm,clip]{HeartQR_153914.png}\\

\url{https://kubo-yuge.github.io/Heart/Heart-copy220503json.html}\\

\columnbreak

\section{かみりん？}

赤い点を動かすと，かみりん？の絵が変形します。\par
\begin{layer}{80}{0}
\putnotese{18}{4}{\scalebox{.45}{%%% /Users/kumac21/DataKuboMac21/予定とメモなど/予定とメモ（学内）/予定：学内22/学事予定22/オープンキャンパス221103/チャレンジキャンパス20のデータ/fig/kamirinOspSample201003.tex 
%%% Generator=kamirinOspSample201003.cdy 
{\unitlength=1cm%
\begin{picture}%
(10,10)(-5,-5)%
\linethickness{0.008in}%%
{%
\color[cmyk]{1,0,0,0.1}%
\polygon*(3.00000,-4.00000)(3.12739,-3.88690)(3.27774,-3.76211)(3.44265,-3.62276)%
(3.61372,-3.46601)(3.78258,-3.28900)(3.94084,-3.08887)(4.08009,-2.86276)(4.19197,-2.60781)%
(4.26806,-2.32118)(4.30000,-2.00000)(4.29263,-1.47057)(4.25090,-0.95034)(4.17867,-0.43914)%
(4.07979,0.06320)(3.95814,0.55684)(3.81758,1.04197)(3.66197,1.51875)(3.49518,1.98735)%
(3.32106,2.44795)(3.14349,2.90071)(3.06116,3.07234)(2.96285,3.21943)(2.85190,3.34439)%
(2.73161,3.44961)(2.60532,3.53746)(2.47634,3.61034)(2.34800,3.67065)(2.22361,3.72077)%
(2.10650,3.76309)(2.00000,3.80000)(1.82325,3.85478)(1.64046,3.89852)(1.45196,3.93245)%
(1.25808,3.95783)(1.05913,3.97589)(0.85546,3.98789)(0.64738,3.99505)(0.43522,3.99863)%
(0.21932,3.99986)(0.00000,4.00000)(-0.21658,3.99968)(-0.43026,3.99799)(-0.64075,3.99379)%
(-0.84774,3.98596)(-1.05096,3.97339)(-1.25009,3.95495)(-1.44486,3.92951)(-1.63496,3.89595)%
(-1.82011,3.85316)(-2.00000,3.80000)(-2.11560,3.76378)(-2.24391,3.72455)(-2.38095,3.67914)%
(-2.52276,3.62439)(-2.66537,3.55713)(-2.80483,3.47419)(-2.93716,3.37241)(-3.05841,3.24863)%
(-3.16460,3.09968)(-3.25177,2.92239)(-3.42683,2.46968)(-3.59450,2.00796)(-3.75171,1.53732)%
(-3.89535,1.05784)(-4.02236,0.56961)(-4.12966,0.07270)(-4.21415,-0.43279)(-4.27276,-0.94680)%
(-4.30240,-1.46923)(-4.30000,-2.00000)(-4.26375,-2.32099)(-4.18507,-2.60758)(-4.07205,-2.86260)%
(-3.93281,-3.08887)(-3.77542,-3.28918)(-3.60801,-3.46636)(-3.43866,-3.62322)(-3.27548,-3.76257)%
(-3.12656,-3.88722)(-3.00000,-4.00000)(-2.93562,-4.05102)(-2.86140,-4.09165)(-2.77805,-4.12322)%
(-2.68631,-4.14702)(-2.58689,-4.16436)(-2.48052,-4.17656)(-2.36793,-4.18491)(-2.24983,-4.19073)%
(-2.12694,-4.19533)(-2.00000,-4.20000)(-1.60012,-4.21439)(-1.20015,-4.22559)(-0.80013,-4.23359)%
(-0.40008,-4.23838)(0.00000,-4.23998)(0.40008,-4.23838)(0.80013,-4.23359)(1.20015,-4.22559)%
(1.60012,-4.21439)(2.00000,-4.20000)(2.12694,-4.19533)(2.24983,-4.19073)(2.36793,-4.18491)%
(2.48052,-4.17656)(2.58689,-4.16436)(2.68631,-4.14702)(2.77805,-4.12322)(2.86140,-4.09165)%
(2.93562,-4.05102)(3.00000,-4.00000)(3.00000,-4.00000)}%
\linethickness{0.016in}%%
\polyline(3.00000,-4.00000)(3.12739,-3.88690)(3.27774,-3.76211)(3.44265,-3.62276)%
(3.61372,-3.46601)(3.78258,-3.28900)(3.94084,-3.08887)(4.08009,-2.86276)(4.19197,-2.60781)%
(4.26806,-2.32118)(4.30000,-2.00000)(4.29263,-1.47057)(4.25090,-0.95034)(4.17867,-0.43914)%
(4.07979,0.06320)(3.95814,0.55684)(3.81758,1.04197)(3.66197,1.51875)(3.49518,1.98735)%
(3.32106,2.44795)(3.14349,2.90071)(3.06116,3.07234)(2.96285,3.21943)(2.85190,3.34439)%
(2.73161,3.44961)(2.60532,3.53746)(2.47634,3.61034)(2.34800,3.67065)(2.22361,3.72077)%
(2.10650,3.76309)(2.00000,3.80000)(1.82325,3.85478)(1.64046,3.89852)(1.45196,3.93245)%
(1.25808,3.95783)(1.05913,3.97589)(0.85546,3.98789)(0.64738,3.99505)(0.43522,3.99863)%
(0.21932,3.99986)(0.00000,4.00000)(-0.21658,3.99968)(-0.43026,3.99799)(-0.64075,3.99379)%
(-0.84774,3.98596)(-1.05096,3.97339)(-1.25009,3.95495)(-1.44486,3.92951)(-1.63496,3.89595)%
(-1.82011,3.85316)(-2.00000,3.80000)(-2.11560,3.76378)(-2.24391,3.72455)(-2.38095,3.67914)%
(-2.52276,3.62439)(-2.66537,3.55713)(-2.80483,3.47419)(-2.93716,3.37241)(-3.05841,3.24863)%
(-3.16460,3.09968)(-3.25177,2.92239)(-3.42683,2.46968)(-3.59450,2.00796)(-3.75171,1.53732)%
(-3.89535,1.05784)(-4.02236,0.56961)(-4.12966,0.07270)(-4.21415,-0.43279)(-4.27276,-0.94680)%
(-4.30240,-1.46923)(-4.30000,-2.00000)(-4.26375,-2.32099)(-4.18507,-2.60758)(-4.07205,-2.86260)%
(-3.93281,-3.08887)(-3.77542,-3.28918)(-3.60801,-3.46636)(-3.43866,-3.62322)(-3.27548,-3.76257)%
(-3.12656,-3.88722)(-3.00000,-4.00000)(-2.93562,-4.05102)(-2.86140,-4.09165)(-2.77805,-4.12322)%
(-2.68631,-4.14702)(-2.58689,-4.16436)(-2.48052,-4.17656)(-2.36793,-4.18491)(-2.24983,-4.19073)%
(-2.12694,-4.19533)(-2.00000,-4.20000)(-1.60012,-4.21439)(-1.20015,-4.22559)(-0.80013,-4.23359)%
(-0.40008,-4.23838)(0.00000,-4.23998)(0.40008,-4.23838)(0.80013,-4.23359)(1.20015,-4.22559)%
(1.60012,-4.21439)(2.00000,-4.20000)(2.12694,-4.19533)(2.24983,-4.19073)(2.36793,-4.18491)%
(2.48052,-4.17656)(2.58689,-4.16436)(2.68631,-4.14702)(2.77805,-4.12322)(2.86140,-4.09165)%
(2.93562,-4.05102)(3.00000,-4.00000)%
%
\linethickness{0.008in}%%
{%
\color[cmyk]{0,0.05,0.2,0}%
\polygon*(3.00000,-4.00000)(3.05432,-3.98532)(3.08381,-3.95330)(3.09326,-3.90834)%
(3.08745,-3.85483)(3.07117,-3.79715)(3.04919,-3.73971)(3.02630,-3.68688)(3.00729,-3.64306)%
(2.99693,-3.61263)(3.00000,-3.60000)(3.09255,-3.58905)(3.18639,-3.59137)(3.28166,-3.60458)%
(3.37848,-3.62626)(3.47701,-3.65402)(3.57738,-3.68546)(3.67972,-3.71818)(3.78418,-3.74977)%
(3.89090,-3.77785)(4.00000,-3.80000)(4.02780,-3.80507)(4.05463,-3.81052)(4.08020,-3.81579)%
(4.10423,-3.82030)(4.12643,-3.82349)(4.14651,-3.82481)(4.16420,-3.82368)(4.17920,-3.81954)%
(4.19123,-3.81184)(4.20000,-3.80000)(4.22117,-3.76105)(4.24690,-3.71647)(4.27455,-3.66662)%
(4.30147,-3.61185)(4.32504,-3.55253)(4.34260,-3.48900)(4.35152,-3.42163)(4.34915,-3.35076)%
(4.33286,-3.27677)(4.30000,-3.20000)(4.24851,-3.11746)(4.18725,-3.04065)(4.11750,-2.96845)%
(4.04057,-2.89977)(3.95772,-2.83351)(3.87026,-2.76857)(3.77947,-2.70384)(3.68664,-2.63824)%
(3.59305,-2.57066)(3.50000,-2.50000)(3.42730,-2.45064)(3.34926,-2.40950)(3.26894,-2.37426)%
(3.18941,-2.34264)(3.11373,-2.31232)(3.04494,-2.28103)(2.98612,-2.24645)(2.94031,-2.20628)%
(2.91059,-2.15823)(2.90000,-2.10000)(2.91607,-2.04654)(2.95979,-1.99074)(3.02446,-1.93044)%
(3.10337,-1.86347)(3.18980,-1.78765)(3.27705,-1.70082)(3.35841,-1.60082)(3.42716,-1.48545)%
(3.47659,-1.35257)(3.50000,-1.20000)(3.51384,-0.91269)(3.52479,-0.61719)(3.52946,-0.31469)%
(3.52443,-0.00637)(3.50630,0.30661)(3.47167,0.62306)(3.41713,0.94180)(3.33928,1.26165)%
(3.23470,1.58145)(3.10000,1.90000)(2.97033,2.13694)(2.82231,2.34659)(2.65916,2.53103)%
(2.48413,2.69235)(2.30046,2.83262)(2.11139,2.95392)(1.92015,3.05832)(1.72998,3.14790)%
(1.54413,3.22475)(1.36583,3.29093)(1.22542,3.32325)(1.08579,3.32462)(0.94696,3.30153)%
(0.80897,3.26046)(0.67184,3.20791)(0.53560,3.15037)(0.40027,3.09432)(0.26587,3.04624)%
(0.13244,3.01264)(0.00000,3.00000)(-0.14667,3.01311)(-0.29522,3.04791)(-0.44553,3.09760)%
(-0.59751,3.15537)(-0.75107,3.21443)(-0.90611,3.26797)(-1.06254,3.30920)(-1.22025,3.33130)%
(-1.37915,3.32748)(-1.53914,3.29093)(-1.71870,3.22347)(-1.90080,3.14256)(-2.08339,3.04728)%
(-2.26436,2.93669)(-2.44165,2.80988)(-2.61318,2.66592)(-2.77685,2.50387)(-2.93060,2.32282)%
(-3.07235,2.12184)(-3.20000,1.90000)(-3.33208,1.61998)(-3.44215,1.33743)(-3.53160,1.05331)%
(-3.60187,0.76857)(-3.65438,0.48414)(-3.69053,0.20099)(-3.71175,-0.07995)(-3.71945,-0.35773)%
(-3.71507,-0.63139)(-3.70000,-0.90000)(-3.66606,-1.09242)(-3.60038,-1.26118)(-3.51104,-1.40913)%
(-3.40609,-1.53909)(-3.29359,-1.65390)(-3.18160,-1.75639)(-3.07819,-1.84942)(-2.99141,-1.93580)%
(-2.92933,-2.01838)(-2.90000,-2.10000)(-2.90371,-2.16560)(-2.93040,-2.21868)(-2.97615,-2.26180)%
(-3.03703,-2.29752)(-3.10910,-2.32841)(-3.18845,-2.35702)(-3.27114,-2.38592)(-3.35324,-2.41768)%
(-3.43084,-2.45485)(-3.50000,-2.50000)(-3.59305,-2.57066)(-3.68664,-2.63824)(-3.77947,-2.70384)%
(-3.87026,-2.76857)(-3.95772,-2.83351)(-4.04057,-2.89977)(-4.11750,-2.96845)(-4.18725,-3.04065)%
(-4.24851,-3.11746)(-4.30000,-3.20000)(-4.33291,-3.27263)(-4.35398,-3.34467)(-4.36408,-3.41528)%
(-4.36406,-3.48363)(-4.35479,-3.54890)(-4.33714,-3.61026)(-4.31197,-3.66688)(-4.28015,-3.71793)%
(-4.24254,-3.76258)(-4.20000,-3.80000)(-4.14486,-3.83317)(-4.08508,-3.85476)(-4.02116,-3.86634)%
(-3.95360,-3.86950)(-3.88291,-3.86581)(-3.80958,-3.85686)(-3.73413,-3.84423)(-3.65704,-3.82948)%
(-3.57883,-3.81422)(-3.50000,-3.80000)(-3.42804,-3.78315)(-3.36197,-3.75919)(-3.30140,-3.73054)%
(-3.24597,-3.69959)(-3.19529,-3.66874)(-3.14899,-3.64039)(-3.10669,-3.61694)(-3.06801,-3.60079)%
(-3.03257,-3.59435)(-3.00000,-3.60000)(-3.00166,-3.61275)(-3.01573,-3.64317)(-3.03741,-3.68692)%
(-3.06192,-3.73962)(-3.08447,-3.79693)(-3.10025,-3.85448)(-3.10449,-3.90792)(-3.09239,-3.95290)%
(-3.05916,-3.98504)(-3.00000,-4.00000)(-2.40048,-4.03596)(-1.80064,-4.06393)(-1.20056,-4.08391)%
(-0.60032,-4.09589)(0.00000,-4.09989)(0.60032,-4.09589)(1.20056,-4.08391)(1.80064,-4.06393)%
(2.40048,-4.03596)(3.00000,-4.00000)(3.00000,-4.00000)}%
\linethickness{0.016in}%%
\polyline(3.00000,-4.00000)(3.05432,-3.98532)(3.08381,-3.95330)(3.09326,-3.90834)%
(3.08745,-3.85483)(3.07117,-3.79715)(3.04919,-3.73971)(3.02630,-3.68688)(3.00729,-3.64306)%
(2.99693,-3.61263)(3.00000,-3.60000)(3.09255,-3.58905)(3.18639,-3.59137)(3.28166,-3.60458)%
(3.37848,-3.62626)(3.47701,-3.65402)(3.57738,-3.68546)(3.67972,-3.71818)(3.78418,-3.74977)%
(3.89090,-3.77785)(4.00000,-3.80000)(4.02780,-3.80507)(4.05463,-3.81052)(4.08020,-3.81579)%
(4.10423,-3.82030)(4.12643,-3.82349)(4.14651,-3.82481)(4.16420,-3.82368)(4.17920,-3.81954)%
(4.19123,-3.81184)(4.20000,-3.80000)(4.22117,-3.76105)(4.24690,-3.71647)(4.27455,-3.66662)%
(4.30147,-3.61185)(4.32504,-3.55253)(4.34260,-3.48900)(4.35152,-3.42163)(4.34915,-3.35076)%
(4.33286,-3.27677)(4.30000,-3.20000)(4.24851,-3.11746)(4.18725,-3.04065)(4.11750,-2.96845)%
(4.04057,-2.89977)(3.95772,-2.83351)(3.87026,-2.76857)(3.77947,-2.70384)(3.68664,-2.63824)%
(3.59305,-2.57066)(3.50000,-2.50000)(3.42730,-2.45064)(3.34926,-2.40950)(3.26894,-2.37426)%
(3.18941,-2.34264)(3.11373,-2.31232)(3.04494,-2.28103)(2.98612,-2.24645)(2.94031,-2.20628)%
(2.91059,-2.15823)(2.90000,-2.10000)(2.91607,-2.04654)(2.95979,-1.99074)(3.02446,-1.93044)%
(3.10337,-1.86347)(3.18980,-1.78765)(3.27705,-1.70082)(3.35841,-1.60082)(3.42716,-1.48545)%
(3.47659,-1.35257)(3.50000,-1.20000)(3.51384,-0.91269)(3.52479,-0.61719)(3.52946,-0.31469)%
(3.52443,-0.00637)(3.50630,0.30661)(3.47167,0.62306)(3.41713,0.94180)(3.33928,1.26165)%
(3.23470,1.58145)(3.10000,1.90000)(2.97033,2.13694)(2.82231,2.34659)(2.65916,2.53103)%
(2.48413,2.69235)(2.30046,2.83262)(2.11139,2.95392)(1.92015,3.05832)(1.72998,3.14790)%
(1.54413,3.22475)(1.36583,3.29093)(1.22542,3.32325)(1.08579,3.32462)(0.94696,3.30153)%
(0.80897,3.26046)(0.67184,3.20791)(0.53560,3.15037)(0.40027,3.09432)(0.26587,3.04624)%
(0.13244,3.01264)(0.00000,3.00000)(-0.14667,3.01311)(-0.29522,3.04791)(-0.44553,3.09760)%
(-0.59751,3.15537)(-0.75107,3.21443)(-0.90611,3.26797)(-1.06254,3.30920)(-1.22025,3.33130)%
(-1.37915,3.32748)(-1.53914,3.29093)(-1.71870,3.22347)(-1.90080,3.14256)(-2.08339,3.04728)%
(-2.26436,2.93669)(-2.44165,2.80988)(-2.61318,2.66592)(-2.77685,2.50387)(-2.93060,2.32282)%
(-3.07235,2.12184)(-3.20000,1.90000)(-3.33208,1.61998)(-3.44215,1.33743)(-3.53160,1.05331)%
(-3.60187,0.76857)(-3.65438,0.48414)(-3.69053,0.20099)(-3.71175,-0.07995)(-3.71945,-0.35773)%
(-3.71507,-0.63139)(-3.70000,-0.90000)(-3.66606,-1.09242)(-3.60038,-1.26118)(-3.51104,-1.40913)%
(-3.40609,-1.53909)(-3.29359,-1.65390)(-3.18160,-1.75639)(-3.07819,-1.84942)(-2.99141,-1.93580)%
(-2.92933,-2.01838)(-2.90000,-2.10000)(-2.90371,-2.16560)(-2.93040,-2.21868)(-2.97615,-2.26180)%
(-3.03703,-2.29752)(-3.10910,-2.32841)(-3.18845,-2.35702)(-3.27114,-2.38592)(-3.35324,-2.41768)%
(-3.43084,-2.45485)(-3.50000,-2.50000)(-3.59305,-2.57066)(-3.68664,-2.63824)(-3.77947,-2.70384)%
(-3.87026,-2.76857)(-3.95772,-2.83351)(-4.04057,-2.89977)(-4.11750,-2.96845)(-4.18725,-3.04065)%
(-4.24851,-3.11746)(-4.30000,-3.20000)(-4.33291,-3.27263)(-4.35398,-3.34467)(-4.36408,-3.41528)%
(-4.36406,-3.48363)(-4.35479,-3.54890)(-4.33714,-3.61026)(-4.31197,-3.66688)(-4.28015,-3.71793)%
(-4.24254,-3.76258)(-4.20000,-3.80000)(-4.14486,-3.83317)(-4.08508,-3.85476)(-4.02116,-3.86634)%
(-3.95360,-3.86950)(-3.88291,-3.86581)(-3.80958,-3.85686)(-3.73413,-3.84423)(-3.65704,-3.82948)%
(-3.57883,-3.81422)(-3.50000,-3.80000)(-3.42804,-3.78315)(-3.36197,-3.75919)(-3.30140,-3.73054)%
(-3.24597,-3.69959)(-3.19529,-3.66874)(-3.14899,-3.64039)(-3.10669,-3.61694)(-3.06801,-3.60079)%
(-3.03257,-3.59435)(-3.00000,-3.60000)(-3.00166,-3.61275)(-3.01573,-3.64317)(-3.03741,-3.68692)%
(-3.06192,-3.73962)(-3.08447,-3.79693)(-3.10025,-3.85448)(-3.10449,-3.90792)(-3.09239,-3.95290)%
(-3.05916,-3.98504)(-3.00000,-4.00000)(-2.40048,-4.03596)(-1.80064,-4.06393)(-1.20056,-4.08391)%
(-0.60032,-4.09589)(0.00000,-4.09989)(0.60032,-4.09589)(1.20056,-4.08391)(1.80064,-4.06393)%
(2.40048,-4.03596)(3.00000,-4.00000)%
%
\linethickness{0.008in}%%
{%
\color[cmyk]{1,0,0,0}%
\polygon*(-0.10000,0.90000)(-0.12904,0.89653)(-0.15327,0.88677)(-0.17290,0.87169)%
(-0.18816,0.85230)(-0.19925,0.82955)(-0.20641,0.80444)(-0.20984,0.77795)(-0.20977,0.75106)%
(-0.20642,0.72475)(-0.20000,0.70000)(-0.19130,0.68045)(-0.17920,0.66312)(-0.16404,0.64799)%
(-0.14615,0.63501)(-0.12586,0.62414)(-0.10349,0.61534)(-0.07940,0.60857)(-0.05389,0.60378)%
(-0.02732,0.60094)(0.00000,0.60000)(0.03184,0.60067)(0.06367,0.60283)(0.09450,0.60671)%
(0.12336,0.61250)(0.14925,0.62045)(0.17121,0.63076)(0.18824,0.64365)(0.19937,0.65934)%
(0.20362,0.67805)(0.20000,0.70000)(0.18660,0.73262)(0.16880,0.76349)(0.14687,0.79223)%
(0.12108,0.81843)(0.09169,0.84169)(0.05896,0.86162)(0.02317,0.87781)(-0.01544,0.88987)%
(-0.05658,0.89740)(-0.10000,0.90000)(-0.10000,0.90000)}%
\linethickness{0.032in}%%
\polyline(-0.10000,0.90000)(-0.12904,0.89653)(-0.15327,0.88677)(-0.17290,0.87169)%
(-0.18816,0.85230)(-0.19925,0.82955)(-0.20641,0.80444)(-0.20984,0.77795)(-0.20977,0.75106)%
(-0.20642,0.72475)(-0.20000,0.70000)(-0.19130,0.68045)(-0.17920,0.66312)(-0.16404,0.64799)%
(-0.14615,0.63501)(-0.12586,0.62414)(-0.10349,0.61534)(-0.07940,0.60857)(-0.05389,0.60378)%
(-0.02732,0.60094)(0.00000,0.60000)(0.03184,0.60067)(0.06367,0.60283)(0.09450,0.60671)%
(0.12336,0.61250)(0.14925,0.62045)(0.17121,0.63076)(0.18824,0.64365)(0.19937,0.65934)%
(0.20362,0.67805)(0.20000,0.70000)(0.18660,0.73262)(0.16880,0.76349)(0.14687,0.79223)%
(0.12108,0.81843)(0.09169,0.84169)(0.05896,0.86162)(0.02317,0.87781)(-0.01544,0.88987)%
(-0.05658,0.89740)(-0.10000,0.90000)%
%
\linethickness{0.008in}%%
{%
\color[cmyk]{0,1,1,0}%
\polygon*(0.60000,0.20000)(0.55776,0.19916)(0.51222,0.18905)(0.46316,0.17189)(0.41030,0.14987)%
(0.35343,0.12520)(0.29228,0.10008)(0.22661,0.07670)(0.15618,0.05728)(0.08075,0.04401)%
(0.00006,0.03910)(-0.08106,0.04469)(-0.15699,0.05971)(-0.22791,0.08148)(-0.29402,0.10736)%
(-0.35550,0.13468)(-0.41254,0.16080)(-0.46532,0.18304)(-0.51403,0.19876)(-0.55886,0.20530)%
(-0.60000,0.20000)(-0.62641,0.18696)(-0.64528,0.16706)(-0.65695,0.14116)(-0.66173,0.11016)%
(-0.65993,0.07491)(-0.65189,0.03629)(-0.63791,-0.00481)(-0.61832,-0.04753)(-0.59344,-0.09100)%
(-0.56359,-0.13433)(-0.52007,-0.19506)(-0.47595,-0.25952)(-0.43039,-0.32513)(-0.38254,-0.38932)%
(-0.33155,-0.44949)(-0.27658,-0.50308)(-0.21679,-0.54751)(-0.15132,-0.58019)(-0.07934,-0.59855)%
(0.00000,-0.60000)(0.08611,-0.58267)(0.16257,-0.54962)(0.23070,-0.50341)(0.29183,-0.44659)%
(0.34727,-0.38172)(0.39834,-0.31134)(0.44637,-0.23802)(0.49267,-0.16431)(0.53857,-0.09276)%
(0.58539,-0.02593)(0.60632,0.00486)(0.62331,0.03579)(0.63620,0.06611)(0.64486,0.09511)%
(0.64914,0.12204)(0.64891,0.14616)(0.64403,0.16676)(0.63434,0.18309)(0.61971,0.19441)%
(0.60000,0.20000)(0.60000,0.20000)}%
\linethickness{0.016in}%%
\polyline(0.60000,0.20000)(0.55776,0.19916)(0.51222,0.18905)(0.46316,0.17189)(0.41030,0.14987)%
(0.35343,0.12520)(0.29228,0.10008)(0.22661,0.07670)(0.15618,0.05728)(0.08075,0.04401)%
(0.00006,0.03910)(-0.08106,0.04469)(-0.15699,0.05971)(-0.22791,0.08148)(-0.29402,0.10736)%
(-0.35550,0.13468)(-0.41254,0.16080)(-0.46532,0.18304)(-0.51403,0.19876)(-0.55886,0.20530)%
(-0.60000,0.20000)(-0.62641,0.18696)(-0.64528,0.16706)(-0.65695,0.14116)(-0.66173,0.11016)%
(-0.65993,0.07491)(-0.65189,0.03629)(-0.63791,-0.00481)(-0.61832,-0.04753)(-0.59344,-0.09100)%
(-0.56359,-0.13433)(-0.52007,-0.19506)(-0.47595,-0.25952)(-0.43039,-0.32513)(-0.38254,-0.38932)%
(-0.33155,-0.44949)(-0.27658,-0.50308)(-0.21679,-0.54751)(-0.15132,-0.58019)(-0.07934,-0.59855)%
(0.00000,-0.60000)(0.08611,-0.58267)(0.16257,-0.54962)(0.23070,-0.50341)(0.29183,-0.44659)%
(0.34727,-0.38172)(0.39834,-0.31134)(0.44637,-0.23802)(0.49267,-0.16431)(0.53857,-0.09276)%
(0.58539,-0.02593)(0.60632,0.00486)(0.62331,0.03579)(0.63620,0.06611)(0.64486,0.09511)%
(0.64914,0.12204)(0.64891,0.14616)(0.64403,0.16676)(0.63434,0.18309)(0.61971,0.19441)%
(0.60000,0.20000)%
%
\linethickness{0.008in}%%
\polygon*(-1.80000,1.70000)(-1.87889,1.69329)(-1.95310,1.67392)(-2.02198,1.64303)%
(-2.08487,1.60177)(-2.14113,1.55126)(-2.19010,1.49265)(-2.23113,1.42708)(-2.26355,1.35569)%
(-2.28673,1.27962)(-2.30000,1.20000)(-2.30187,1.12408)(-2.29265,1.05015)(-2.27300,0.97962)%
(-2.24360,0.91392)(-2.20512,0.85444)(-2.15822,0.80261)(-2.10359,0.75984)(-2.04190,0.72754)%
(-1.97381,0.70712)(-1.90000,0.70000)(-1.81153,0.70595)(-1.72594,0.72337)(-1.64451,0.75163)%
(-1.56855,0.79010)(-1.49934,0.83816)(-1.43819,0.89516)(-1.38638,0.96047)(-1.34522,1.03347)%
(-1.31599,1.11353)(-1.30000,1.20000)(-1.29956,1.28238)(-1.31418,1.36006)(-1.34240,1.43210)%
(-1.38280,1.49757)(-1.43392,1.55553)(-1.49433,1.60505)(-1.56258,1.64519)(-1.63724,1.67502)%
(-1.71686,1.69360)(-1.80000,1.70000)(1.70000,1.70000)(1.78110,1.69346)(1.85804,1.67451)%
(1.92978,1.64419)(1.99529,1.60353)(2.05355,1.55355)(2.10353,1.49529)(2.14419,1.42978)%
(2.17451,1.35804)(2.19346,1.28110)(2.20000,1.20000)(2.19346,1.11890)(2.17451,1.04196)%
(2.14419,0.97022)(2.10353,0.90471)(2.05355,0.84645)(1.99529,0.79647)(1.92978,0.75581)%
(1.85804,0.72549)(1.78110,0.70654)(1.70000,0.70000)(1.61890,0.70654)(1.54196,0.72549)%
(1.47022,0.75581)(1.40471,0.79647)(1.34645,0.84645)(1.29647,0.90471)(1.25581,0.97022)%
(1.22549,1.04196)(1.20654,1.11890)(1.20000,1.20000)(1.20654,1.28110)(1.22549,1.35804)%
(1.25581,1.42978)(1.29647,1.49529)(1.34645,1.55355)(1.40471,1.60353)(1.47022,1.64419)%
(1.54196,1.67451)(1.61890,1.69346)(1.70000,1.70000)(-1.80000,1.70000)\linethickness{0.016in}%%
\polyline(-1.80000,1.70000)(-1.87889,1.69329)(-1.95310,1.67392)(-2.02198,1.64303)%
(-2.08487,1.60177)(-2.14113,1.55126)(-2.19010,1.49265)(-2.23113,1.42708)(-2.26355,1.35569)%
(-2.28673,1.27962)(-2.30000,1.20000)(-2.30187,1.12408)(-2.29265,1.05015)(-2.27300,0.97962)%
(-2.24360,0.91392)(-2.20512,0.85444)(-2.15822,0.80261)(-2.10359,0.75984)(-2.04190,0.72754)%
(-1.97381,0.70712)(-1.90000,0.70000)(-1.81153,0.70595)(-1.72594,0.72337)(-1.64451,0.75163)%
(-1.56855,0.79010)(-1.49934,0.83816)(-1.43819,0.89516)(-1.38638,0.96047)(-1.34522,1.03347)%
(-1.31599,1.11353)(-1.30000,1.20000)(-1.29956,1.28238)(-1.31418,1.36006)(-1.34240,1.43210)%
(-1.38280,1.49757)(-1.43392,1.55553)(-1.49433,1.60505)(-1.56258,1.64519)(-1.63724,1.67502)%
(-1.71686,1.69360)(-1.80000,1.70000)%
%
\linethickness{0.008in}%%
\linethickness{0.016in}%%
\polyline(1.70000,1.70000)(1.61890,1.69346)(1.54196,1.67451)(1.47022,1.64419)(1.40471,1.60353)%
(1.34645,1.55355)(1.29647,1.49529)(1.25581,1.42978)(1.22549,1.35804)(1.20654,1.28110)%
(1.20000,1.20000)(1.20654,1.11890)(1.22549,1.04196)(1.25581,0.97022)(1.29647,0.90471)%
(1.34645,0.84645)(1.40471,0.79647)(1.47022,0.75581)(1.54196,0.72549)(1.61890,0.70654)%
(1.70000,0.70000)(1.78110,0.70654)(1.85804,0.72549)(1.92978,0.75581)(1.99529,0.79647)%
(2.05355,0.84645)(2.10353,0.90471)(2.14419,0.97022)(2.17451,1.04196)(2.19346,1.11890)%
(2.20000,1.20000)(2.19346,1.28110)(2.17451,1.35804)(2.14419,1.42978)(2.10353,1.49529)%
(2.05355,1.55355)(1.99529,1.60353)(1.92978,1.64419)(1.85804,1.67451)(1.78110,1.69346)%
(1.70000,1.70000)%
%
\linethickness{0.008in}%%
{%
\color[cmyk]{0,0.8,0.9,0}%
{%
\color[cmyk]{0,0.8,0.9,0}%
\polygon*(-2.80000,-2.30000)(-2.80158,-2.27493)(-2.80628,-2.25026)(-2.81404,-2.22638)%
(-2.82474,-2.20365)(-2.83820,-2.18244)(-2.85421,-2.16309)(-2.87252,-2.14590)(-2.89283,-2.13113)%
(-2.91484,-2.11903)(-2.93820,-2.10979)(-2.96252,-2.10354)(-2.98744,-2.10039)(-3.01256,-2.10039)%
(-3.03748,-2.10354)(-3.06180,-2.10979)(-3.08516,-2.11903)(-3.10717,-2.13113)(-3.12748,-2.14590)%
(-3.14579,-2.16309)(-3.16180,-2.18244)(-3.17526,-2.20365)(-3.18596,-2.22638)(-3.19372,-2.25026)%
(-3.19842,-2.27493)(-3.20000,-2.30000)(-3.19842,-2.32507)(-3.19372,-2.34974)(-3.18596,-2.37362)%
(-3.17526,-2.39635)(-3.16180,-2.41756)(-3.14579,-2.43691)(-3.12748,-2.45410)(-3.10717,-2.46887)%
(-3.08516,-2.48097)(-3.06180,-2.49021)(-3.03748,-2.49646)(-3.01256,-2.49961)(-2.98744,-2.49961)%
(-2.96252,-2.49646)(-2.93820,-2.49021)(-2.91484,-2.48097)(-2.89283,-2.46887)(-2.87252,-2.45410)%
(-2.85421,-2.43691)(-2.83820,-2.41756)(-2.82474,-2.39635)(-2.81404,-2.37362)(-2.80628,-2.34974)%
(-2.80158,-2.32507)(-2.80000,-2.30000)(-2.40000,-3.30000)(-2.40158,-3.32507)(-2.40628,-3.34974)%
(-2.41404,-3.37362)(-2.42474,-3.39635)(-2.43820,-3.41756)(-2.45421,-3.43691)(-2.47252,-3.45410)%
(-2.49283,-3.46887)(-2.51484,-3.48097)(-2.53820,-3.49021)(-2.56252,-3.49646)(-2.58744,-3.49961)%
(-2.61256,-3.49961)(-2.63748,-3.49646)(-2.66180,-3.49021)(-2.68516,-3.48097)(-2.70717,-3.46887)%
(-2.72748,-3.45410)(-2.74579,-3.43691)(-2.76180,-3.41756)(-2.77526,-3.39635)(-2.78596,-3.37362)%
(-2.79372,-3.34974)(-2.79842,-3.32507)(-2.80000,-3.30000)(-2.79842,-3.27493)(-2.79372,-3.25026)%
(-2.78596,-3.22638)(-2.77526,-3.20365)(-2.76180,-3.18244)(-2.74579,-3.16309)(-2.72748,-3.14590)%
(-2.70717,-3.13113)(-2.68516,-3.11903)(-2.66180,-3.10979)(-2.63748,-3.10354)(-2.61256,-3.10039)%
(-2.58744,-3.10039)(-2.56252,-3.10354)(-2.53820,-3.10979)(-2.51484,-3.11903)(-2.49283,-3.13113)%
(-2.47252,-3.14590)(-2.45421,-3.16309)(-2.43820,-3.18244)(-2.42474,-3.20365)(-2.41404,-3.22638)%
(-2.40628,-3.25026)(-2.40158,-3.27493)(-2.40000,-3.30000)(-2.80000,-2.30000)}%
\polyline(-2.80000,-2.30000)(-2.80158,-2.27493)(-2.80628,-2.25026)(-2.81404,-2.22638)%
(-2.82474,-2.20365)(-2.83820,-2.18244)(-2.85421,-2.16309)(-2.87252,-2.14590)(-2.89283,-2.13113)%
(-2.91484,-2.11903)(-2.93820,-2.10979)(-2.96252,-2.10354)(-2.98744,-2.10039)(-3.01256,-2.10039)%
(-3.03748,-2.10354)(-3.06180,-2.10979)(-3.08516,-2.11903)(-3.10717,-2.13113)(-3.12748,-2.14590)%
(-3.14579,-2.16309)(-3.16180,-2.18244)(-3.17526,-2.20365)(-3.18596,-2.22638)(-3.19372,-2.25026)%
(-3.19842,-2.27493)(-3.20000,-2.30000)(-3.19842,-2.32507)(-3.19372,-2.34974)(-3.18596,-2.37362)%
(-3.17526,-2.39635)(-3.16180,-2.41756)(-3.14579,-2.43691)(-3.12748,-2.45410)(-3.10717,-2.46887)%
(-3.08516,-2.48097)(-3.06180,-2.49021)(-3.03748,-2.49646)(-3.01256,-2.49961)(-2.98744,-2.49961)%
(-2.96252,-2.49646)(-2.93820,-2.49021)(-2.91484,-2.48097)(-2.89283,-2.46887)(-2.87252,-2.45410)%
(-2.85421,-2.43691)(-2.83820,-2.41756)(-2.82474,-2.39635)(-2.81404,-2.37362)(-2.80628,-2.34974)%
(-2.80158,-2.32507)(-2.80000,-2.30000)%
%
}%
{%
\color[cmyk]{0,0.3,0.5,0}%
{%
\color[cmyk]{0,0.3,0.6,0}%
\polygon*(-2.70000,-2.80000)(-2.70158,-2.77493)(-2.70628,-2.75026)(-2.71404,-2.72638)%
(-2.72474,-2.70365)(-2.73820,-2.68244)(-2.75421,-2.66309)(-2.77252,-2.64590)(-2.79283,-2.63113)%
(-2.81484,-2.61903)(-2.83820,-2.60979)(-2.86252,-2.60354)(-2.88744,-2.60039)(-2.91256,-2.60039)%
(-2.93748,-2.60354)(-2.96180,-2.60979)(-2.98516,-2.61903)(-3.00717,-2.63113)(-3.02748,-2.64590)%
(-3.04579,-2.66309)(-3.06180,-2.68244)(-3.07526,-2.70365)(-3.08596,-2.72638)(-3.09372,-2.75026)%
(-3.09842,-2.77493)(-3.10000,-2.80000)(-3.09842,-2.82507)(-3.09372,-2.84974)(-3.08596,-2.87362)%
(-3.07526,-2.89635)(-3.06180,-2.91756)(-3.04579,-2.93691)(-3.02748,-2.95410)(-3.00717,-2.96887)%
(-2.98516,-2.98097)(-2.96180,-2.99021)(-2.93748,-2.99646)(-2.91256,-2.99961)(-2.88744,-2.99961)%
(-2.86252,-2.99646)(-2.83820,-2.99021)(-2.81484,-2.98097)(-2.79283,-2.96887)(-2.77252,-2.95410)%
(-2.75421,-2.93691)(-2.73820,-2.91756)(-2.72474,-2.89635)(-2.71404,-2.87362)(-2.70628,-2.84974)%
(-2.70158,-2.82507)(-2.70000,-2.80000)(-2.00000,-3.70000)(-2.00158,-3.72507)(-2.00628,-3.74974)%
(-2.01404,-3.77362)(-2.02474,-3.79635)(-2.03820,-3.81756)(-2.05421,-3.83691)(-2.07252,-3.85410)%
(-2.09283,-3.86887)(-2.11484,-3.88097)(-2.13820,-3.89021)(-2.16252,-3.89646)(-2.18744,-3.89961)%
(-2.21256,-3.89961)(-2.23748,-3.89646)(-2.26180,-3.89021)(-2.28516,-3.88097)(-2.30717,-3.86887)%
(-2.32748,-3.85410)(-2.34579,-3.83691)(-2.36180,-3.81756)(-2.37526,-3.79635)(-2.38596,-3.77362)%
(-2.39372,-3.74974)(-2.39842,-3.72507)(-2.40000,-3.70000)(-2.39842,-3.67493)(-2.39372,-3.65026)%
(-2.38596,-3.62638)(-2.37526,-3.60365)(-2.36180,-3.58244)(-2.34579,-3.56309)(-2.32748,-3.54590)%
(-2.30717,-3.53113)(-2.28516,-3.51903)(-2.26180,-3.50979)(-2.23748,-3.50354)(-2.21256,-3.50039)%
(-2.18744,-3.50039)(-2.16252,-3.50354)(-2.13820,-3.50979)(-2.11484,-3.51903)(-2.09283,-3.53113)%
(-2.07252,-3.54590)(-2.05421,-3.56309)(-2.03820,-3.58244)(-2.02474,-3.60365)(-2.01404,-3.62638)%
(-2.00628,-3.65026)(-2.00158,-3.67493)(-2.00000,-3.70000)(-2.70000,-2.80000)}%
\polyline(-2.70000,-2.80000)(-2.70158,-2.77493)(-2.70628,-2.75026)(-2.71404,-2.72638)%
(-2.72474,-2.70365)(-2.73820,-2.68244)(-2.75421,-2.66309)(-2.77252,-2.64590)(-2.79283,-2.63113)%
(-2.81484,-2.61903)(-2.83820,-2.60979)(-2.86252,-2.60354)(-2.88744,-2.60039)(-2.91256,-2.60039)%
(-2.93748,-2.60354)(-2.96180,-2.60979)(-2.98516,-2.61903)(-3.00717,-2.63113)(-3.02748,-2.64590)%
(-3.04579,-2.66309)(-3.06180,-2.68244)(-3.07526,-2.70365)(-3.08596,-2.72638)(-3.09372,-2.75026)%
(-3.09842,-2.77493)(-3.10000,-2.80000)(-3.09842,-2.82507)(-3.09372,-2.84974)(-3.08596,-2.87362)%
(-3.07526,-2.89635)(-3.06180,-2.91756)(-3.04579,-2.93691)(-3.02748,-2.95410)(-3.00717,-2.96887)%
(-2.98516,-2.98097)(-2.96180,-2.99021)(-2.93748,-2.99646)(-2.91256,-2.99961)(-2.88744,-2.99961)%
(-2.86252,-2.99646)(-2.83820,-2.99021)(-2.81484,-2.98097)(-2.79283,-2.96887)(-2.77252,-2.95410)%
(-2.75421,-2.93691)(-2.73820,-2.91756)(-2.72474,-2.89635)(-2.71404,-2.87362)(-2.70628,-2.84974)%
(-2.70158,-2.82507)(-2.70000,-2.80000)%
%
}%
{%
\color[cmyk]{0,0.8,0.9,0}%
\polyline(-2.40000,-3.30000)(-2.40158,-3.27493)(-2.40628,-3.25026)(-2.41404,-3.22638)%
(-2.42474,-3.20365)(-2.43820,-3.18244)(-2.45421,-3.16309)(-2.47252,-3.14590)(-2.49283,-3.13113)%
(-2.51484,-3.11903)(-2.53820,-3.10979)(-2.56252,-3.10354)(-2.58744,-3.10039)(-2.61256,-3.10039)%
(-2.63748,-3.10354)(-2.66180,-3.10979)(-2.68516,-3.11903)(-2.70717,-3.13113)(-2.72748,-3.14590)%
(-2.74579,-3.16309)(-2.76180,-3.18244)(-2.77526,-3.20365)(-2.78596,-3.22638)(-2.79372,-3.25026)%
(-2.79842,-3.27493)(-2.80000,-3.30000)(-2.79842,-3.32507)(-2.79372,-3.34974)(-2.78596,-3.37362)%
(-2.77526,-3.39635)(-2.76180,-3.41756)(-2.74579,-3.43691)(-2.72748,-3.45410)(-2.70717,-3.46887)%
(-2.68516,-3.48097)(-2.66180,-3.49021)(-2.63748,-3.49646)(-2.61256,-3.49961)(-2.58744,-3.49961)%
(-2.56252,-3.49646)(-2.53820,-3.49021)(-2.51484,-3.48097)(-2.49283,-3.46887)(-2.47252,-3.45410)%
(-2.45421,-3.43691)(-2.43820,-3.41756)(-2.42474,-3.39635)(-2.41404,-3.37362)(-2.40628,-3.34974)%
(-2.40158,-3.32507)(-2.40000,-3.30000)%
%
}%
{%
\color[cmyk]{0,0.3,0.5,0}%
\polyline(-2.00000,-3.70000)(-2.00158,-3.67493)(-2.00628,-3.65026)(-2.01404,-3.62638)%
(-2.02474,-3.60365)(-2.03820,-3.58244)(-2.05421,-3.56309)(-2.07252,-3.54590)(-2.09283,-3.53113)%
(-2.11484,-3.51903)(-2.13820,-3.50979)(-2.16252,-3.50354)(-2.18744,-3.50039)(-2.21256,-3.50039)%
(-2.23748,-3.50354)(-2.26180,-3.50979)(-2.28516,-3.51903)(-2.30717,-3.53113)(-2.32748,-3.54590)%
(-2.34579,-3.56309)(-2.36180,-3.58244)(-2.37526,-3.60365)(-2.38596,-3.62638)(-2.39372,-3.65026)%
(-2.39842,-3.67493)(-2.40000,-3.70000)(-2.39842,-3.72507)(-2.39372,-3.74974)(-2.38596,-3.77362)%
(-2.37526,-3.79635)(-2.36180,-3.81756)(-2.34579,-3.83691)(-2.32748,-3.85410)(-2.30717,-3.86887)%
(-2.28516,-3.88097)(-2.26180,-3.89021)(-2.23748,-3.89646)(-2.21256,-3.89961)(-2.18744,-3.89961)%
(-2.16252,-3.89646)(-2.13820,-3.89021)(-2.11484,-3.88097)(-2.09283,-3.86887)(-2.07252,-3.85410)%
(-2.05421,-3.83691)(-2.03820,-3.81756)(-2.02474,-3.79635)(-2.01404,-3.77362)(-2.00628,-3.74974)%
(-2.00158,-3.72507)(-2.00000,-3.70000)%
%
}%
{%
\color[cmyk]{0,0.2,0.2,0}%
{%
\color[cmyk]{0,0.2,0.2,0}%
\polygon*(-2.00000,0.00000)(-2.00394,0.06267)(-2.01571,0.12434)(-2.03511,0.18406)%
(-2.06185,0.24088)(-2.09549,0.29389)(-2.13552,0.34227)(-2.18129,0.38526)(-2.23209,0.42216)%
(-2.28711,0.45241)(-2.34549,0.47553)(-2.40631,0.49114)(-2.46860,0.49901)(-2.53140,0.49901)%
(-2.59369,0.49114)(-2.65451,0.47553)(-2.71289,0.45241)(-2.76791,0.42216)(-2.81871,0.38526)%
(-2.86448,0.34227)(-2.90451,0.29389)(-2.93815,0.24088)(-2.96489,0.18406)(-2.98429,0.12434)%
(-2.99606,0.06267)(-3.00000,-0.00000)(-2.99606,-0.06267)(-2.98429,-0.12434)(-2.96489,-0.18406)%
(-2.93815,-0.24088)(-2.90451,-0.29389)(-2.86448,-0.34227)(-2.81871,-0.38526)(-2.76791,-0.42216)%
(-2.71289,-0.45241)(-2.65451,-0.47553)(-2.59369,-0.49114)(-2.53140,-0.49901)(-2.46860,-0.49901)%
(-2.40631,-0.49114)(-2.34549,-0.47553)(-2.28711,-0.45241)(-2.23209,-0.42216)(-2.18129,-0.38526)%
(-2.13552,-0.34227)(-2.09549,-0.29389)(-2.06185,-0.24088)(-2.03511,-0.18406)(-2.01571,-0.12434)%
(-2.00394,-0.06267)(-2.00000,-0.00000)(2.80000,-0.00000)(2.79606,-0.06267)(2.78429,-0.12434)%
(2.76489,-0.18406)(2.73815,-0.24088)(2.70451,-0.29389)(2.66448,-0.34227)(2.61871,-0.38526)%
(2.56791,-0.42216)(2.51289,-0.45241)(2.45451,-0.47553)(2.39369,-0.49114)(2.33140,-0.49901)%
(2.26860,-0.49901)(2.20631,-0.49114)(2.14549,-0.47553)(2.08711,-0.45241)(2.03209,-0.42216)%
(1.98129,-0.38526)(1.93552,-0.34227)(1.89549,-0.29389)(1.86185,-0.24088)(1.83511,-0.18406)%
(1.81571,-0.12434)(1.80394,-0.06267)(1.80000,-0.00000)(1.80394,0.06267)(1.81571,0.12434)%
(1.83511,0.18406)(1.86185,0.24088)(1.89549,0.29389)(1.93552,0.34227)(1.98129,0.38526)%
(2.03209,0.42216)(2.08711,0.45241)(2.14549,0.47553)(2.20631,0.49114)(2.26860,0.49901)%
(2.33140,0.49901)(2.39369,0.49114)(2.45451,0.47553)(2.51289,0.45241)(2.56791,0.42216)%
(2.61871,0.38526)(2.66448,0.34227)(2.70451,0.29389)(2.73815,0.24088)(2.76489,0.18406)%
(2.78429,0.12434)(2.79606,0.06267)(2.80000,0.00000)(-2.00000,0.00000)}%
\polyline(-2.00000,0.00000)(-2.00394,0.06267)(-2.01571,0.12434)(-2.03511,0.18406)%
(-2.06185,0.24088)(-2.09549,0.29389)(-2.13552,0.34227)(-2.18129,0.38526)(-2.23209,0.42216)%
(-2.28711,0.45241)(-2.34549,0.47553)(-2.40631,0.49114)(-2.46860,0.49901)(-2.53140,0.49901)%
(-2.59369,0.49114)(-2.65451,0.47553)(-2.71289,0.45241)(-2.76791,0.42216)(-2.81871,0.38526)%
(-2.86448,0.34227)(-2.90451,0.29389)(-2.93815,0.24088)(-2.96489,0.18406)(-2.98429,0.12434)%
(-2.99606,0.06267)(-3.00000,-0.00000)(-2.99606,-0.06267)(-2.98429,-0.12434)(-2.96489,-0.18406)%
(-2.93815,-0.24088)(-2.90451,-0.29389)(-2.86448,-0.34227)(-2.81871,-0.38526)(-2.76791,-0.42216)%
(-2.71289,-0.45241)(-2.65451,-0.47553)(-2.59369,-0.49114)(-2.53140,-0.49901)(-2.46860,-0.49901)%
(-2.40631,-0.49114)(-2.34549,-0.47553)(-2.28711,-0.45241)(-2.23209,-0.42216)(-2.18129,-0.38526)%
(-2.13552,-0.34227)(-2.09549,-0.29389)(-2.06185,-0.24088)(-2.03511,-0.18406)(-2.01571,-0.12434)%
(-2.00394,-0.06267)(-2.00000,-0.00000)%
%
}%
{%
\color[cmyk]{0,0.4,0.4,0}%
{%
\color[cmyk]{0,0.4,0.4,0}%
\polygon*(-2.20000,0.00000)(-2.20237,0.03760)(-2.20943,0.07461)(-2.22107,0.11044)%
(-2.23711,0.14453)(-2.25729,0.17634)(-2.28131,0.20536)(-2.30877,0.23115)(-2.33925,0.25330)%
(-2.37227,0.27145)(-2.40729,0.28532)(-2.44379,0.29469)(-2.48116,0.29941)(-2.51884,0.29941)%
(-2.55621,0.29469)(-2.59271,0.28532)(-2.62773,0.27145)(-2.66075,0.25330)(-2.69123,0.23115)%
(-2.71869,0.20536)(-2.74271,0.17634)(-2.76289,0.14453)(-2.77893,0.11044)(-2.79057,0.07461)%
(-2.79763,0.03760)(-2.80000,-0.00000)(-2.79763,-0.03760)(-2.79057,-0.07461)(-2.77893,-0.11044)%
(-2.76289,-0.14453)(-2.74271,-0.17634)(-2.71869,-0.20536)(-2.69123,-0.23115)(-2.66075,-0.25330)%
(-2.62773,-0.27145)(-2.59271,-0.28532)(-2.55621,-0.29469)(-2.51884,-0.29941)(-2.48116,-0.29941)%
(-2.44379,-0.29469)(-2.40729,-0.28532)(-2.37227,-0.27145)(-2.33925,-0.25330)(-2.30877,-0.23115)%
(-2.28131,-0.20536)(-2.25729,-0.17634)(-2.23711,-0.14453)(-2.22107,-0.11044)(-2.20943,-0.07461)%
(-2.20237,-0.03760)(-2.20000,-0.00000)(2.60000,-0.00000)(2.59763,-0.03760)(2.59057,-0.07461)%
(2.57893,-0.11044)(2.56289,-0.14453)(2.54271,-0.17634)(2.51869,-0.20536)(2.49123,-0.23115)%
(2.46075,-0.25330)(2.42773,-0.27145)(2.39271,-0.28532)(2.35621,-0.29469)(2.31884,-0.29941)%
(2.28116,-0.29941)(2.24379,-0.29469)(2.20729,-0.28532)(2.17227,-0.27145)(2.13925,-0.25330)%
(2.10877,-0.23115)(2.08131,-0.20536)(2.05729,-0.17634)(2.03711,-0.14453)(2.02107,-0.11044)%
(2.00943,-0.07461)(2.00237,-0.03760)(2.00000,-0.00000)(2.00237,0.03760)(2.00943,0.07461)%
(2.02107,0.11044)(2.03711,0.14453)(2.05729,0.17634)(2.08131,0.20536)(2.10877,0.23115)%
(2.13925,0.25330)(2.17227,0.27145)(2.20729,0.28532)(2.24379,0.29469)(2.28116,0.29941)%
(2.31884,0.29941)(2.35621,0.29469)(2.39271,0.28532)(2.42773,0.27145)(2.46075,0.25330)%
(2.49123,0.23115)(2.51869,0.20536)(2.54271,0.17634)(2.56289,0.14453)(2.57893,0.11044)%
(2.59057,0.07461)(2.59763,0.03760)(2.60000,0.00000)(-2.20000,0.00000)}%
\polyline(-2.20000,0.00000)(-2.20237,0.03760)(-2.20943,0.07461)(-2.22107,0.11044)%
(-2.23711,0.14453)(-2.25729,0.17634)(-2.28131,0.20536)(-2.30877,0.23115)(-2.33925,0.25330)%
(-2.37227,0.27145)(-2.40729,0.28532)(-2.44379,0.29469)(-2.48116,0.29941)(-2.51884,0.29941)%
(-2.55621,0.29469)(-2.59271,0.28532)(-2.62773,0.27145)(-2.66075,0.25330)(-2.69123,0.23115)%
(-2.71869,0.20536)(-2.74271,0.17634)(-2.76289,0.14453)(-2.77893,0.11044)(-2.79057,0.07461)%
(-2.79763,0.03760)(-2.80000,-0.00000)(-2.79763,-0.03760)(-2.79057,-0.07461)(-2.77893,-0.11044)%
(-2.76289,-0.14453)(-2.74271,-0.17634)(-2.71869,-0.20536)(-2.69123,-0.23115)(-2.66075,-0.25330)%
(-2.62773,-0.27145)(-2.59271,-0.28532)(-2.55621,-0.29469)(-2.51884,-0.29941)(-2.48116,-0.29941)%
(-2.44379,-0.29469)(-2.40729,-0.28532)(-2.37227,-0.27145)(-2.33925,-0.25330)(-2.30877,-0.23115)%
(-2.28131,-0.20536)(-2.25729,-0.17634)(-2.23711,-0.14453)(-2.22107,-0.11044)(-2.20943,-0.07461)%
(-2.20237,-0.03760)(-2.20000,-0.00000)%
%
}%
{%
\color[cmyk]{0,0.2,0.2,0}%
\polyline(2.80000,0.00000)(2.79606,0.06267)(2.78429,0.12434)(2.76489,0.18406)(2.73815,0.24088)%
(2.70451,0.29389)(2.66448,0.34227)(2.61871,0.38526)(2.56791,0.42216)(2.51289,0.45241)%
(2.45451,0.47553)(2.39369,0.49114)(2.33140,0.49901)(2.26860,0.49901)(2.20631,0.49114)%
(2.14549,0.47553)(2.08711,0.45241)(2.03209,0.42216)(1.98129,0.38526)(1.93552,0.34227)%
(1.89549,0.29389)(1.86185,0.24088)(1.83511,0.18406)(1.81571,0.12434)(1.80394,0.06267)%
(1.80000,-0.00000)(1.80394,-0.06267)(1.81571,-0.12434)(1.83511,-0.18406)(1.86185,-0.24088)%
(1.89549,-0.29389)(1.93552,-0.34227)(1.98129,-0.38526)(2.03209,-0.42216)(2.08711,-0.45241)%
(2.14549,-0.47553)(2.20631,-0.49114)(2.26860,-0.49901)(2.33140,-0.49901)(2.39369,-0.49114)%
(2.45451,-0.47553)(2.51289,-0.45241)(2.56791,-0.42216)(2.61871,-0.38526)(2.66448,-0.34227)%
(2.70451,-0.29389)(2.73815,-0.24088)(2.76489,-0.18406)(2.78429,-0.12434)(2.79606,-0.06267)%
(2.80000,-0.00000)%
%
}%
{%
\color[cmyk]{0,0.4,0.4,0}%
\polyline(2.60000,0.00000)(2.59763,0.03760)(2.59057,0.07461)(2.57893,0.11044)(2.56289,0.14453)%
(2.54271,0.17634)(2.51869,0.20536)(2.49123,0.23115)(2.46075,0.25330)(2.42773,0.27145)%
(2.39271,0.28532)(2.35621,0.29469)(2.31884,0.29941)(2.28116,0.29941)(2.24379,0.29469)%
(2.20729,0.28532)(2.17227,0.27145)(2.13925,0.25330)(2.10877,0.23115)(2.08131,0.20536)%
(2.05729,0.17634)(2.03711,0.14453)(2.02107,0.11044)(2.00943,0.07461)(2.00237,0.03760)%
(2.00000,-0.00000)(2.00237,-0.03760)(2.00943,-0.07461)(2.02107,-0.11044)(2.03711,-0.14453)%
(2.05729,-0.17634)(2.08131,-0.20536)(2.10877,-0.23115)(2.13925,-0.25330)(2.17227,-0.27145)%
(2.20729,-0.28532)(2.24379,-0.29469)(2.28116,-0.29941)(2.31884,-0.29941)(2.35621,-0.29469)%
(2.39271,-0.28532)(2.42773,-0.27145)(2.46075,-0.25330)(2.49123,-0.23115)(2.51869,-0.20536)%
(2.54271,-0.17634)(2.56289,-0.14453)(2.57893,-0.11044)(2.59057,-0.07461)(2.59763,-0.03760)%
(2.60000,-0.00000)%
%
}%
{%
\color[cmyk]{1,0,0.4,0.3}%
\polygon*(-3.50000,1.00000)(-3.44319,1.05241)(-3.39726,1.10662)(-3.36121,1.16177)%
(-3.33401,1.21697)(-3.31465,1.27134)(-3.30212,1.32400)(-3.29540,1.37408)(-3.29349,1.42069)%
(-3.29536,1.46296)(-3.30000,1.50000)(-3.31232,1.52328)(-3.33740,1.53674)(-3.37192,1.54296)%
(-3.41257,1.54450)(-3.45605,1.54395)(-3.49905,1.54388)(-3.53825,1.54688)(-3.57035,1.55551)%
(-3.59204,1.57236)(-3.60000,1.60000)(-3.59074,1.66492)(-3.56575,1.72655)(-3.52919,1.78573)%
(-3.48524,1.84328)(-3.43807,1.90000)(-3.39185,1.95672)(-3.35077,2.01427)(-3.31898,2.07345)%
(-3.30067,2.13508)(-3.30000,2.20000)(-3.31445,2.23370)(-3.34353,2.25280)(-3.38387,2.26081)%
(-3.43213,2.26122)(-3.48494,2.25753)(-3.53896,2.25324)(-3.59083,2.25184)(-3.63720,2.25683)%
(-3.67471,2.27172)(-3.70000,2.30000)(-3.71335,2.34840)(-3.70957,2.39942)(-3.69314,2.45222)%
(-3.66854,2.50595)(-3.64025,2.55975)(-3.61274,2.61277)(-3.59048,2.66415)(-3.57795,2.71306)%
(-3.57963,2.75862)(-3.60000,2.80000)(-3.63450,2.83465)(-3.67463,2.86033)(-3.71955,2.87754)%
(-3.76846,2.88673)(-3.82054,2.88838)(-3.87498,2.88296)(-3.93096,2.87094)(-3.98767,2.85279)%
(-4.04428,2.82899)(-4.10000,2.80000)(-4.14035,2.77255)(-4.17460,2.74047)(-4.20335,2.70445)%
(-4.22719,2.66518)(-4.24670,2.62335)(-4.26248,2.57966)(-4.27511,2.53478)(-4.28518,2.48942)%
(-4.29328,2.44426)(-4.30000,2.40000)(-4.30018,2.36733)(-4.29263,2.33549)(-4.27939,2.30440)%
(-4.26253,2.27396)(-4.24409,2.24409)(-4.22612,2.21468)(-4.21067,2.18566)(-4.19980,2.15694)%
(-4.19556,2.12841)(-4.20000,2.10000)(-4.21785,2.06072)(-4.24240,2.02623)(-4.27215,1.99550)%
(-4.30560,1.96751)(-4.34125,1.94125)(-4.37760,1.91568)(-4.41315,1.88980)(-4.44640,1.86257)%
(-4.47585,1.83298)(-4.50000,1.80000)(-4.51175,1.77678)(-4.51939,1.75384)(-4.52349,1.73135)%
(-4.52461,1.70946)(-4.52331,1.68835)(-4.52014,1.66816)(-4.51566,1.64907)(-4.51044,1.63124)%
(-4.50503,1.61483)(-4.50000,1.60000)(-4.49407,1.59031)(-4.48403,1.58180)(-4.47109,1.57401)%
(-4.45644,1.56647)(-4.44129,1.55871)(-4.42684,1.55024)(-4.41429,1.54061)(-4.40483,1.52934)%
(-4.39967,1.51596)(-4.40000,1.50000)(-4.41279,1.45684)(-4.43216,1.41578)(-4.45625,1.37631)%
(-4.48318,1.33789)(-4.51110,1.30000)(-4.53812,1.26211)(-4.56239,1.22369)(-4.58204,1.18422)%
(-4.59520,1.14316)(-4.60000,1.10000)(-4.59518,1.08156)(-4.58198,1.06874)(-4.56228,1.06015)%
(-4.53794,1.05437)(-4.51086,1.05000)(-4.48291,1.04563)(-4.45598,1.03985)(-4.43192,1.03126)%
(-4.41264,1.01844)(-4.40000,1.00000)(-4.39507,0.96416)(-4.40195,0.91946)(-4.41661,0.86888)%
(-4.43505,0.81539)(-4.45324,0.76197)(-4.46717,0.71160)(-4.47283,0.66723)(-4.46620,0.63187)%
(-4.44326,0.60846)(-4.40000,0.60000)(-4.31974,0.60412)(-4.23648,0.61643)(-4.15054,0.63684)%
(-4.06222,0.66529)(-3.97185,0.70167)(-3.87973,0.74593)(-3.78617,0.79797)(-3.69149,0.85771)%
(-3.59599,0.92508)(-3.50000,1.00000)(3.30000,0.80000)(3.38192,0.70962)(3.46256,0.62296)%
(3.54224,0.54116)(3.62128,0.46534)(3.70000,0.39666)(3.77872,0.33624)(3.85776,0.28523)%
(3.93744,0.24476)(4.01808,0.21597)(4.10000,0.20000)(4.15270,0.20440)(4.18546,0.22925)%
(4.20221,0.27068)(4.20685,0.32482)(4.20331,0.38780)(4.19551,0.45576)(4.18736,0.52483)%
(4.18278,0.59114)(4.18568,0.65081)(4.20000,0.70000)(4.21383,0.71799)(4.23318,0.73065)%
(4.25649,0.73932)(4.28221,0.74533)(4.30879,0.75000)(4.33466,0.75467)(4.35828,0.76068)%
(4.37807,0.76935)(4.39250,0.78201)(4.40000,0.80000)(4.40324,0.85021)(4.39765,0.90028)%
(4.38544,0.95025)(4.36882,1.00014)(4.35000,1.05000)(4.33118,1.09986)(4.31456,1.14975)%
(4.30235,1.19972)(4.29676,1.24979)(4.30000,1.30000)(4.30887,1.32063)(4.32597,1.33498)%
(4.34914,1.34452)(4.37622,1.35072)(4.40505,1.35505)(4.43348,1.35898)(4.45935,1.36397)%
(4.48049,1.37149)(4.49476,1.38301)(4.50000,1.40000)(4.49440,1.43190)(4.47920,1.46409)%
(4.45680,1.49713)(4.42960,1.53160)(4.40000,1.56806)(4.37040,1.60707)(4.34320,1.64920)%
(4.32080,1.69502)(4.30560,1.74510)(4.30000,1.80000)(4.30758,1.84913)(4.32784,1.89774)%
(4.35705,1.94610)(4.39151,1.99447)(4.42749,2.04313)(4.46126,2.09234)(4.48912,2.14236)%
(4.50735,2.19347)(4.51221,2.24592)(4.50000,2.30000)(4.45761,2.39429)(4.40937,2.48061)%
(4.35599,2.55828)(4.29821,2.62657)(4.23676,2.68479)(4.17237,2.73222)(4.10577,2.76817)%
(4.03769,2.79192)(3.96885,2.80276)(3.90000,2.80000)(3.84991,2.78080)(3.82120,2.74330)%
(3.80920,2.69097)(3.80923,2.62727)(3.81660,2.55566)(3.82665,2.47960)(3.83469,2.40254)%
(3.83605,2.32795)(3.82604,2.25928)(3.80000,2.20000)(3.77600,2.17590)(3.74695,2.16471)%
(3.71425,2.16361)(3.67928,2.16975)(3.64344,2.18031)(3.60812,2.19244)(3.57472,2.20331)%
(3.54464,2.21008)(3.51927,2.20992)(3.50000,2.20000)(3.48589,2.18152)(3.47671,2.15984)%
(3.47178,2.13520)(3.47044,2.10784)(3.47201,2.07799)(3.47581,2.04591)(3.48119,2.01183)%
(3.48746,1.97599)(3.49395,1.93863)(3.50000,1.90000)(3.51004,1.84998)(3.52444,1.80226)%
(3.54156,1.75680)(3.55976,1.71359)(3.57741,1.67259)(3.59287,1.63378)(3.60450,1.59715)%
(3.61065,1.56265)(3.60970,1.53028)(3.60000,1.50000)(3.58351,1.48093)(3.56058,1.47329)%
(3.53250,1.47438)(3.50056,1.48151)(3.46605,1.49197)(3.43025,1.50308)(3.39446,1.51213)%
(3.35996,1.51644)(3.32805,1.51329)(3.30000,1.50000)(3.27917,1.48158)(3.26111,1.45990)%
(3.24568,1.43520)(3.23275,1.40775)(3.22218,1.37782)(3.21383,1.34566)(3.20758,1.31154)%
(3.20328,1.27571)(3.20080,1.23845)(3.20000,1.20000)(3.19942,1.16891)(3.19839,1.13438)%
(3.19796,1.09696)(3.19917,1.05720)(3.20309,1.01564)(3.21076,0.97283)(3.22323,0.92931)%
(3.24156,0.88564)(3.26680,0.84235)(3.30000,0.80000)(-3.50000,1.00000)}%
\linethickness{0.016in}%%
\polyline(-3.50000,1.00000)(-3.44319,1.05241)(-3.39726,1.10662)(-3.36121,1.16177)%
(-3.33401,1.21697)(-3.31465,1.27134)(-3.30212,1.32400)(-3.29540,1.37408)(-3.29349,1.42069)%
(-3.29536,1.46296)(-3.30000,1.50000)(-3.31232,1.52328)(-3.33740,1.53674)(-3.37192,1.54296)%
(-3.41257,1.54450)(-3.45605,1.54395)(-3.49905,1.54388)(-3.53825,1.54688)(-3.57035,1.55551)%
(-3.59204,1.57236)(-3.60000,1.60000)(-3.59074,1.66492)(-3.56575,1.72655)(-3.52919,1.78573)%
(-3.48524,1.84328)(-3.43807,1.90000)(-3.39185,1.95672)(-3.35077,2.01427)(-3.31898,2.07345)%
(-3.30067,2.13508)(-3.30000,2.20000)(-3.31445,2.23370)(-3.34353,2.25280)(-3.38387,2.26081)%
(-3.43213,2.26122)(-3.48494,2.25753)(-3.53896,2.25324)(-3.59083,2.25184)(-3.63720,2.25683)%
(-3.67471,2.27172)(-3.70000,2.30000)(-3.71335,2.34840)(-3.70957,2.39942)(-3.69314,2.45222)%
(-3.66854,2.50595)(-3.64025,2.55975)(-3.61274,2.61277)(-3.59048,2.66415)(-3.57795,2.71306)%
(-3.57963,2.75862)(-3.60000,2.80000)(-3.63450,2.83465)(-3.67463,2.86033)(-3.71955,2.87754)%
(-3.76846,2.88673)(-3.82054,2.88838)(-3.87498,2.88296)(-3.93096,2.87094)(-3.98767,2.85279)%
(-4.04428,2.82899)(-4.10000,2.80000)(-4.14035,2.77255)(-4.17460,2.74047)(-4.20335,2.70445)%
(-4.22719,2.66518)(-4.24670,2.62335)(-4.26248,2.57966)(-4.27511,2.53478)(-4.28518,2.48942)%
(-4.29328,2.44426)(-4.30000,2.40000)(-4.30018,2.36733)(-4.29263,2.33549)(-4.27939,2.30440)%
(-4.26253,2.27396)(-4.24409,2.24409)(-4.22612,2.21468)(-4.21067,2.18566)(-4.19980,2.15694)%
(-4.19556,2.12841)(-4.20000,2.10000)(-4.21785,2.06072)(-4.24240,2.02623)(-4.27215,1.99550)%
(-4.30560,1.96751)(-4.34125,1.94125)(-4.37760,1.91568)(-4.41315,1.88980)(-4.44640,1.86257)%
(-4.47585,1.83298)(-4.50000,1.80000)(-4.51175,1.77678)(-4.51939,1.75384)(-4.52349,1.73135)%
(-4.52461,1.70946)(-4.52331,1.68835)(-4.52014,1.66816)(-4.51566,1.64907)(-4.51044,1.63124)%
(-4.50503,1.61483)(-4.50000,1.60000)(-4.49407,1.59031)(-4.48403,1.58180)(-4.47109,1.57401)%
(-4.45644,1.56647)(-4.44129,1.55871)(-4.42684,1.55024)(-4.41429,1.54061)(-4.40483,1.52934)%
(-4.39967,1.51596)(-4.40000,1.50000)(-4.41279,1.45684)(-4.43216,1.41578)(-4.45625,1.37631)%
(-4.48318,1.33789)(-4.51110,1.30000)(-4.53812,1.26211)(-4.56239,1.22369)(-4.58204,1.18422)%
(-4.59520,1.14316)(-4.60000,1.10000)(-4.59518,1.08156)(-4.58198,1.06874)(-4.56228,1.06015)%
(-4.53794,1.05437)(-4.51086,1.05000)(-4.48291,1.04563)(-4.45598,1.03985)(-4.43192,1.03126)%
(-4.41264,1.01844)(-4.40000,1.00000)(-4.39507,0.96416)(-4.40195,0.91946)(-4.41661,0.86888)%
(-4.43505,0.81539)(-4.45324,0.76197)(-4.46717,0.71160)(-4.47283,0.66723)(-4.46620,0.63187)%
(-4.44326,0.60846)(-4.40000,0.60000)(-4.31974,0.60412)(-4.23648,0.61643)(-4.15054,0.63684)%
(-4.06222,0.66529)(-3.97185,0.70167)(-3.87973,0.74593)(-3.78617,0.79797)(-3.69149,0.85771)%
(-3.59599,0.92508)(-3.50000,1.00000)%
%
\linethickness{0.008in}%%
\linethickness{0.016in}%%
\polyline(3.30000,0.80000)(3.26680,0.84235)(3.24156,0.88564)(3.22323,0.92931)(3.21076,0.97283)%
(3.20309,1.01564)(3.19917,1.05720)(3.19796,1.09696)(3.19839,1.13438)(3.19942,1.16891)%
(3.20000,1.20000)(3.20080,1.23845)(3.20328,1.27571)(3.20758,1.31154)(3.21383,1.34566)%
(3.22218,1.37782)(3.23275,1.40775)(3.24568,1.43520)(3.26111,1.45990)(3.27917,1.48158)%
(3.30000,1.50000)(3.32805,1.51329)(3.35996,1.51644)(3.39446,1.51213)(3.43025,1.50308)%
(3.46605,1.49197)(3.50056,1.48151)(3.53250,1.47438)(3.56058,1.47329)(3.58351,1.48093)%
(3.60000,1.50000)(3.60970,1.53028)(3.61065,1.56265)(3.60450,1.59715)(3.59287,1.63378)%
(3.57741,1.67259)(3.55976,1.71359)(3.54156,1.75680)(3.52444,1.80226)(3.51004,1.84998)%
(3.50000,1.90000)(3.49395,1.93863)(3.48746,1.97599)(3.48119,2.01183)(3.47581,2.04591)%
(3.47201,2.07799)(3.47044,2.10784)(3.47178,2.13520)(3.47671,2.15984)(3.48589,2.18152)%
(3.50000,2.20000)(3.51927,2.20992)(3.54464,2.21008)(3.57472,2.20331)(3.60812,2.19244)%
(3.64344,2.18031)(3.67928,2.16975)(3.71425,2.16361)(3.74695,2.16471)(3.77600,2.17590)%
(3.80000,2.20000)(3.82604,2.25928)(3.83605,2.32795)(3.83469,2.40254)(3.82665,2.47960)%
(3.81660,2.55566)(3.80923,2.62727)(3.80920,2.69097)(3.82120,2.74330)(3.84991,2.78080)%
(3.90000,2.80000)(3.96885,2.80276)(4.03769,2.79192)(4.10577,2.76817)(4.17237,2.73222)%
(4.23676,2.68479)(4.29821,2.62657)(4.35599,2.55828)(4.40937,2.48061)(4.45761,2.39429)%
(4.50000,2.30000)(4.51221,2.24592)(4.50735,2.19347)(4.48912,2.14236)(4.46126,2.09234)%
(4.42749,2.04313)(4.39151,1.99447)(4.35705,1.94610)(4.32784,1.89774)(4.30758,1.84913)%
(4.30000,1.80000)(4.30560,1.74510)(4.32080,1.69502)(4.34320,1.64920)(4.37040,1.60707)%
(4.40000,1.56806)(4.42960,1.53160)(4.45680,1.49713)(4.47920,1.46409)(4.49440,1.43190)%
(4.50000,1.40000)(4.49476,1.38301)(4.48049,1.37149)(4.45935,1.36397)(4.43348,1.35898)%
(4.40505,1.35505)(4.37622,1.35072)(4.34914,1.34452)(4.32597,1.33498)(4.30887,1.32063)%
(4.30000,1.30000)(4.29676,1.24979)(4.30235,1.19972)(4.31456,1.14975)(4.33118,1.09986)%
(4.35000,1.05000)(4.36882,1.00014)(4.38544,0.95025)(4.39765,0.90028)(4.40324,0.85021)%
(4.40000,0.80000)(4.39250,0.78201)(4.37807,0.76935)(4.35828,0.76068)(4.33466,0.75467)%
(4.30879,0.75000)(4.28221,0.74533)(4.25649,0.73932)(4.23318,0.73065)(4.21383,0.71799)%
(4.20000,0.70000)(4.18568,0.65081)(4.18278,0.59114)(4.18736,0.52483)(4.19551,0.45576)%
(4.20331,0.38780)(4.20685,0.32482)(4.20221,0.27068)(4.18546,0.22925)(4.15270,0.20440)%
(4.10000,0.20000)(4.01808,0.21597)(3.93744,0.24476)(3.85776,0.28523)(3.77872,0.33624)%
(3.70000,0.39666)(3.62128,0.46534)(3.54224,0.54116)(3.46256,0.62296)(3.38192,0.70962)%
(3.30000,0.80000)%
%
\linethickness{0.008in}%%
{%
\color[cmyk]{0,0.5,0.2,0}%
\polygon*(-3.20000,0.70000)(-3.19334,0.69405)(-3.19546,0.68773)(-3.20444,0.68085)%
(-3.21839,0.67321)(-3.23541,0.66459)(-3.25359,0.65481)(-3.27104,0.64367)(-3.28586,0.63095)%
(-3.29615,0.61646)(-3.30000,0.60000)(-3.29663,0.58135)(-3.28758,0.56307)(-3.27443,0.54489)%
(-3.25875,0.52658)(-3.24212,0.50788)(-3.22613,0.48855)(-3.21234,0.46834)(-3.20233,0.44702)%
(-3.19770,0.42432)(-3.20000,0.40000)(-3.20773,0.37685)(-3.21891,0.35462)(-3.23327,0.33323)%
(-3.25055,0.31259)(-3.27046,0.29262)(-3.29274,0.27324)(-3.31710,0.25437)(-3.34328,0.23592)%
(-3.37101,0.21783)(-3.40000,0.20000)(-3.44008,0.17920)(-3.47984,0.16400)(-3.51856,0.15320)%
(-3.55552,0.14560)(-3.59000,0.14000)(-3.62128,0.13520)(-3.64864,0.13000)(-3.67136,0.12320)%
(-3.68872,0.11360)(-3.70000,0.10000)(-3.70215,0.08557)(-3.69298,0.07582)(-3.67489,0.06936)%
(-3.65027,0.06480)(-3.62153,0.06077)(-3.59107,0.05587)(-3.56129,0.04873)(-3.53458,0.03796)%
(-3.51335,0.02218)(-3.50000,0.00000)(-3.49485,-0.02217)(-3.49391,-0.04436)(-3.49686,-0.06640)%
(-3.50339,-0.08807)(-3.51318,-0.10920)(-3.52592,-0.12960)(-3.54129,-0.14907)(-3.55897,-0.16742)%
(-3.57864,-0.18446)(-3.60000,-0.20000)(-3.61758,-0.20799)(-3.63678,-0.21065)(-3.65718,-0.20932)%
(-3.67839,-0.20533)(-3.70000,-0.20000)(-3.72161,-0.19467)(-3.74282,-0.19068)(-3.76322,-0.18935)%
(-3.78242,-0.19201)(-3.80000,-0.20000)(-3.82127,-0.21964)(-3.83667,-0.24445)(-3.84746,-0.27251)%
(-3.85490,-0.30193)(-3.86026,-0.33079)(-3.86480,-0.35719)(-3.86978,-0.37921)(-3.87646,-0.39496)%
(-3.88612,-0.40253)(-3.90000,-0.40000)(-3.91470,-0.38140)(-3.92249,-0.34764)(-3.92557,-0.30346)%
(-3.92613,-0.25359)(-3.92639,-0.20277)(-3.92853,-0.15573)(-3.93476,-0.11719)(-3.94728,-0.09191)%
(-3.96830,-0.08460)(-4.00000,-0.10000)(-4.04470,-0.14289)(-4.07960,-0.19467)(-4.10715,-0.25342)%
(-4.12980,-0.31725)(-4.15000,-0.38424)(-4.17020,-0.45250)(-4.19285,-0.52010)(-4.22040,-0.58516)%
(-4.25530,-0.64576)(-4.30000,-0.70000)(-4.31897,-0.71302)(-4.33305,-0.71057)(-4.34333,-0.69568)%
(-4.35088,-0.67142)(-4.35681,-0.64084)(-4.36219,-0.60699)(-4.36811,-0.57293)(-4.37566,-0.54171)%
(-4.38593,-0.51638)(-4.40000,-0.50000)(-4.42190,-0.48997)(-4.44531,-0.48663)(-4.46952,-0.48830)%
(-4.49380,-0.49331)(-4.51741,-0.50000)(-4.53963,-0.50669)(-4.55972,-0.51170)(-4.57697,-0.51337)%
(-4.59064,-0.51003)(-4.60000,-0.50000)(-4.60328,-0.48722)(-4.60138,-0.47331)(-4.59515,-0.45809)%
(-4.58542,-0.44138)(-4.57301,-0.42301)(-4.55877,-0.40280)(-4.54351,-0.38058)(-4.52808,-0.35615)%
(-4.51329,-0.32935)(-4.50000,-0.30000)(-4.48377,-0.26339)(-4.46572,-0.22820)(-4.44705,-0.19445)%
(-4.42897,-0.16216)(-4.41268,-0.13134)(-4.39937,-0.10202)(-4.39025,-0.07420)(-4.38652,-0.04792)%
(-4.38937,-0.02318)(-4.40000,0.00000)(-4.42413,0.02313)(-4.45423,0.03290)(-4.48855,0.03239)%
(-4.52534,0.02467)(-4.56285,0.01285)(-4.59933,0.00000)(-4.63304,-0.01080)(-4.66222,-0.01645)%
(-4.68512,-0.01388)(-4.70000,0.00000)(-4.70266,0.01717)(-4.69389,0.03255)(-4.67608,0.04706)%
(-4.65163,0.06159)(-4.62295,0.07705)(-4.59243,0.09435)(-4.56248,0.11439)(-4.53549,0.13807)%
(-4.51386,0.16631)(-4.50000,0.20000)(-4.49913,0.22291)(-4.50667,0.24505)(-4.52038,0.26646)%
(-4.53800,0.28719)(-4.55729,0.30729)(-4.57599,0.32680)(-4.59187,0.34578)(-4.60266,0.36428)%
(-4.60612,0.38233)(-4.60000,0.40000)(-4.58246,0.42178)(-4.56130,0.44040)(-4.53684,0.45605)%
(-4.50940,0.46895)(-4.47929,0.47929)(-4.44684,0.48730)(-4.41237,0.49316)(-4.37619,0.49710)%
(-4.33863,0.49931)(-4.30000,0.50000)(-4.27655,0.49550)(-4.25352,0.48356)(-4.23102,0.46650)%
(-4.20920,0.44667)(-4.18820,0.42639)(-4.16813,0.40800)(-4.14915,0.39384)(-4.13138,0.38623)%
(-4.11495,0.38750)(-4.10000,0.40000)(-4.07685,0.44029)(-4.06246,0.48517)(-4.05466,0.53372)%
(-4.05123,0.58505)(-4.05000,0.63822)(-4.04877,0.69234)(-4.04534,0.74649)(-4.03754,0.79976)%
(-4.02315,0.85123)(-4.00000,0.90000)(-3.98575,0.91654)(-3.96851,0.92540)(-3.94890,0.92807)%
(-3.92759,0.92607)(-3.90522,0.92089)(-3.88243,0.91404)(-3.85987,0.90702)(-3.83818,0.90134)%
(-3.81801,0.89850)(-3.80000,0.90000)(-3.77651,0.90964)(-3.75301,0.92537)(-3.72979,0.94493)%
(-3.70715,0.96607)(-3.68539,0.98654)(-3.66479,1.00408)(-3.64566,1.01646)(-3.62828,1.02140)%
(-3.61297,1.01667)(-3.60000,1.00000)(-3.57993,0.94705)(-3.57062,0.89344)(-3.56970,0.83967)%
(-3.57482,0.78623)(-3.58361,0.73361)(-3.59370,0.68230)(-3.60275,0.63279)(-3.60839,0.58557)%
(-3.60826,0.54115)(-3.60000,0.50000)(-3.59401,0.48713)(-3.58960,0.48909)(-3.58583,0.50307)%
(-3.58179,0.52623)(-3.57654,0.55576)(-3.56918,0.58883)(-3.55876,0.62261)(-3.54438,0.65428)%
(-3.52510,0.68102)(-3.50000,0.70000)(-3.46271,0.71615)(-3.42650,0.72720)(-3.39157,0.73371)%
(-3.35814,0.73626)(-3.32639,0.73541)(-3.29653,0.73173)(-3.26877,0.72578)(-3.24329,0.71813)%
(-3.22030,0.70935)(-3.20000,0.70000)(2.90000,0.80000)(2.91818,0.80622)(2.93922,0.80832)%
(2.96305,0.80647)(2.98956,0.80090)(3.01867,0.79178)(3.05028,0.77932)(3.08431,0.76371)%
(3.12067,0.74516)(3.15926,0.72386)(3.20000,0.70000)(3.23252,0.67562)(3.26086,0.64581)%
(3.28556,0.61298)(3.30712,0.57952)(3.32608,0.54783)(3.34295,0.52032)(3.35826,0.49938)%
(3.37252,0.48741)(3.38626,0.48682)(3.40000,0.50000)(3.42244,0.54906)(3.42991,0.60186)%
(3.42618,0.65708)(3.41496,0.71339)(3.40000,0.76948)(3.38504,0.82400)(3.37382,0.87564)%
(3.37009,0.92307)(3.37756,0.96497)(3.40000,1.00000)(3.41934,1.01123)(3.43476,1.00370)%
(3.44778,0.98101)(3.45993,0.94675)(3.47273,0.90453)(3.48772,0.85795)(3.50641,0.81061)%
(3.53034,0.76610)(3.56103,0.72803)(3.60000,0.70000)(3.62111,0.69704)(3.64443,0.70567)%
(3.66904,0.72263)(3.69403,0.74465)(3.71846,0.76846)(3.74141,0.79079)(3.76197,0.80838)%
(3.77920,0.81796)(3.79218,0.81625)(3.80000,0.80000)(3.80579,0.74025)(3.80105,0.68034)%
(3.78842,0.62030)(3.77053,0.56017)(3.75000,0.50000)(3.72947,0.43983)(3.71158,0.37970)%
(3.69895,0.31966)(3.69421,0.25975)(3.70000,0.20000)(3.70732,0.18384)(3.71927,0.18129)%
(3.73516,0.18943)(3.75431,0.20533)(3.77606,0.22606)(3.79973,0.24871)(3.82463,0.27036)%
(3.85009,0.28807)(3.87544,0.29892)(3.90000,0.30000)(3.93912,0.29083)(3.97692,0.27920)%
(4.01314,0.26513)(4.04751,0.24865)(4.07977,0.22977)(4.10964,0.20850)(4.13687,0.18487)%
(4.16118,0.15890)(4.18231,0.13060)(4.20000,0.10000)(4.20587,0.07899)(4.20330,0.05866)%
(4.19413,0.03882)(4.18021,0.01933)(4.16338,0.00000)(4.14547,-0.01933)(4.12834,-0.03882)%
(4.11382,-0.05866)(4.10376,-0.07899)(4.10000,-0.10000)(4.10280,-0.12457)(4.11040,-0.14738)%
(4.12160,-0.16869)(4.13520,-0.18880)(4.15000,-0.20799)(4.16480,-0.22653)(4.17840,-0.24473)%
(4.18960,-0.26284)(4.19720,-0.28118)(4.20000,-0.30000)(4.19897,-0.31519)(4.19589,-0.33085)%
(4.19079,-0.34638)(4.18368,-0.36116)(4.17458,-0.37458)(4.16352,-0.38604)(4.15051,-0.39492)%
(4.13557,-0.40061)(4.11873,-0.40251)(4.10000,-0.40000)(4.05933,-0.38241)(4.01962,-0.35297)%
(3.98141,-0.31592)(3.94522,-0.27551)(3.91158,-0.23597)(3.88101,-0.20155)(3.85404,-0.17650)%
(3.83120,-0.16507)(3.81301,-0.17148)(3.80000,-0.20000)(3.79433,-0.25563)(3.80299,-0.31139)%
(3.82274,-0.36759)(3.85036,-0.42456)(3.88260,-0.48260)(3.91624,-0.54204)(3.94803,-0.60318)%
(3.97475,-0.66634)(3.99315,-0.73185)(4.00000,-0.80000)(3.99617,-0.82372)(3.98549,-0.83152)%
(3.96918,-0.82654)(3.94846,-0.81195)(3.92456,-0.79093)(3.89869,-0.76664)(3.87208,-0.74223)%
(3.84595,-0.72088)(3.82151,-0.70575)(3.80000,-0.70000)(3.78980,-0.70434)(3.77974,-0.71588)%
(3.76977,-0.73239)(3.75987,-0.75164)(3.75000,-0.77141)(3.74013,-0.78946)(3.73023,-0.80357)%
(3.72026,-0.81152)(3.71020,-0.81107)(3.70000,-0.80000)(3.66677,-0.73788)(3.63990,-0.67472)%
(3.61811,-0.61108)(3.60012,-0.54753)(3.58465,-0.48465)(3.57040,-0.42299)(3.55610,-0.36313)%
(3.54045,-0.30563)(3.52218,-0.25107)(3.50000,-0.20000)(3.49050,-0.18245)(3.48734,-0.18071)%
(3.48892,-0.19196)(3.49367,-0.21343)(3.50000,-0.24230)(3.50633,-0.27579)(3.51108,-0.31110)%
(3.51266,-0.34544)(3.50950,-0.37600)(3.50000,-0.40000)(3.48446,-0.42136)(3.46742,-0.44103)%
(3.44907,-0.45871)(3.42960,-0.47408)(3.40920,-0.48682)(3.38807,-0.49661)(3.36640,-0.50314)%
(3.34436,-0.50609)(3.32217,-0.50515)(3.30000,-0.50000)(3.28192,-0.49073)(3.26787,-0.47703)%
(3.25701,-0.45965)(3.24850,-0.43938)(3.24151,-0.41698)(3.23520,-0.39322)(3.22873,-0.36887)%
(3.22127,-0.34471)(3.21197,-0.32149)(3.20000,-0.30000)(3.19156,-0.28957)(3.18134,-0.28016)%
(3.16988,-0.27143)(3.15772,-0.26303)(3.14539,-0.25461)(3.13343,-0.24582)(3.12237,-0.23631)%
(3.11276,-0.22574)(3.10512,-0.21375)(3.10000,-0.20000)(3.09197,-0.16751)(3.08443,-0.13506)%
(3.07786,-0.10284)(3.07275,-0.07104)(3.06959,-0.03986)(3.06886,-0.00949)(3.07105,0.01987)%
(3.07664,0.04804)(3.08613,0.07481)(3.10000,0.10000)(3.11650,0.11308)(3.14018,0.11744)%
(3.16863,0.11526)(3.19947,0.10872)(3.23028,0.10000)(3.25867,0.09128)(3.28223,0.08474)%
(3.29858,0.08256)(3.30530,0.08692)(3.30000,0.10000)(3.27449,0.12939)(3.24144,0.15494)%
(3.20247,0.17815)(3.15920,0.20051)(3.11325,0.22350)(3.06624,0.24861)(3.01979,0.27733)%
(2.97552,0.31114)(2.93505,0.35154)(2.90000,0.40000)(2.89154,0.42270)(2.89261,0.44452)%
(2.90117,0.46558)(2.91519,0.48597)(2.93263,0.50579)(2.95146,0.52515)(2.96964,0.54415)%
(2.98515,0.56289)(2.99595,0.58147)(3.00000,0.60000)(2.99580,0.63001)(2.98464,0.65800)%
(2.96863,0.68389)(2.94992,0.70760)(2.93062,0.72906)(2.91288,0.74820)(2.89881,0.76494)%
(2.89056,0.77920)(2.89024,0.79091)(2.90000,0.80000)(-3.20000,0.70000)}%
\linethickness{0.016in}%%
\polyline(-3.20000,0.70000)(-3.19334,0.69405)(-3.19546,0.68773)(-3.20444,0.68085)%
(-3.21839,0.67321)(-3.23541,0.66459)(-3.25359,0.65481)(-3.27104,0.64367)(-3.28586,0.63095)%
(-3.29615,0.61646)(-3.30000,0.60000)(-3.29663,0.58135)(-3.28758,0.56307)(-3.27443,0.54489)%
(-3.25875,0.52658)(-3.24212,0.50788)(-3.22613,0.48855)(-3.21234,0.46834)(-3.20233,0.44702)%
(-3.19770,0.42432)(-3.20000,0.40000)(-3.20773,0.37685)(-3.21891,0.35462)(-3.23327,0.33323)%
(-3.25055,0.31259)(-3.27046,0.29262)(-3.29274,0.27324)(-3.31710,0.25437)(-3.34328,0.23592)%
(-3.37101,0.21783)(-3.40000,0.20000)(-3.44008,0.17920)(-3.47984,0.16400)(-3.51856,0.15320)%
(-3.55552,0.14560)(-3.59000,0.14000)(-3.62128,0.13520)(-3.64864,0.13000)(-3.67136,0.12320)%
(-3.68872,0.11360)(-3.70000,0.10000)(-3.70215,0.08557)(-3.69298,0.07582)(-3.67489,0.06936)%
(-3.65027,0.06480)(-3.62153,0.06077)(-3.59107,0.05587)(-3.56129,0.04873)(-3.53458,0.03796)%
(-3.51335,0.02218)(-3.50000,0.00000)(-3.49485,-0.02217)(-3.49391,-0.04436)(-3.49686,-0.06640)%
(-3.50339,-0.08807)(-3.51318,-0.10920)(-3.52592,-0.12960)(-3.54129,-0.14907)(-3.55897,-0.16742)%
(-3.57864,-0.18446)(-3.60000,-0.20000)(-3.61758,-0.20799)(-3.63678,-0.21065)(-3.65718,-0.20932)%
(-3.67839,-0.20533)(-3.70000,-0.20000)(-3.72161,-0.19467)(-3.74282,-0.19068)(-3.76322,-0.18935)%
(-3.78242,-0.19201)(-3.80000,-0.20000)(-3.82127,-0.21964)(-3.83667,-0.24445)(-3.84746,-0.27251)%
(-3.85490,-0.30193)(-3.86026,-0.33079)(-3.86480,-0.35719)(-3.86978,-0.37921)(-3.87646,-0.39496)%
(-3.88612,-0.40253)(-3.90000,-0.40000)(-3.91470,-0.38140)(-3.92249,-0.34764)(-3.92557,-0.30346)%
(-3.92613,-0.25359)(-3.92639,-0.20277)(-3.92853,-0.15573)(-3.93476,-0.11719)(-3.94728,-0.09191)%
(-3.96830,-0.08460)(-4.00000,-0.10000)(-4.04470,-0.14289)(-4.07960,-0.19467)(-4.10715,-0.25342)%
(-4.12980,-0.31725)(-4.15000,-0.38424)(-4.17020,-0.45250)(-4.19285,-0.52010)(-4.22040,-0.58516)%
(-4.25530,-0.64576)(-4.30000,-0.70000)(-4.31897,-0.71302)(-4.33305,-0.71057)(-4.34333,-0.69568)%
(-4.35088,-0.67142)(-4.35681,-0.64084)(-4.36219,-0.60699)(-4.36811,-0.57293)(-4.37566,-0.54171)%
(-4.38593,-0.51638)(-4.40000,-0.50000)(-4.42190,-0.48997)(-4.44531,-0.48663)(-4.46952,-0.48830)%
(-4.49380,-0.49331)(-4.51741,-0.50000)(-4.53963,-0.50669)(-4.55972,-0.51170)(-4.57697,-0.51337)%
(-4.59064,-0.51003)(-4.60000,-0.50000)(-4.60328,-0.48722)(-4.60138,-0.47331)(-4.59515,-0.45809)%
(-4.58542,-0.44138)(-4.57301,-0.42301)(-4.55877,-0.40280)(-4.54351,-0.38058)(-4.52808,-0.35615)%
(-4.51329,-0.32935)(-4.50000,-0.30000)(-4.48377,-0.26339)(-4.46572,-0.22820)(-4.44705,-0.19445)%
(-4.42897,-0.16216)(-4.41268,-0.13134)(-4.39937,-0.10202)(-4.39025,-0.07420)(-4.38652,-0.04792)%
(-4.38937,-0.02318)(-4.40000,0.00000)(-4.42413,0.02313)(-4.45423,0.03290)(-4.48855,0.03239)%
(-4.52534,0.02467)(-4.56285,0.01285)(-4.59933,0.00000)(-4.63304,-0.01080)(-4.66222,-0.01645)%
(-4.68512,-0.01388)(-4.70000,0.00000)(-4.70266,0.01717)(-4.69389,0.03255)(-4.67608,0.04706)%
(-4.65163,0.06159)(-4.62295,0.07705)(-4.59243,0.09435)(-4.56248,0.11439)(-4.53549,0.13807)%
(-4.51386,0.16631)(-4.50000,0.20000)(-4.49913,0.22291)(-4.50667,0.24505)(-4.52038,0.26646)%
(-4.53800,0.28719)(-4.55729,0.30729)(-4.57599,0.32680)(-4.59187,0.34578)(-4.60266,0.36428)%
(-4.60612,0.38233)(-4.60000,0.40000)(-4.58246,0.42178)(-4.56130,0.44040)(-4.53684,0.45605)%
(-4.50940,0.46895)(-4.47929,0.47929)(-4.44684,0.48730)(-4.41237,0.49316)(-4.37619,0.49710)%
(-4.33863,0.49931)(-4.30000,0.50000)(-4.27655,0.49550)(-4.25352,0.48356)(-4.23102,0.46650)%
(-4.20920,0.44667)(-4.18820,0.42639)(-4.16813,0.40800)(-4.14915,0.39384)(-4.13138,0.38623)%
(-4.11495,0.38750)(-4.10000,0.40000)(-4.07685,0.44029)(-4.06246,0.48517)(-4.05466,0.53372)%
(-4.05123,0.58505)(-4.05000,0.63822)(-4.04877,0.69234)(-4.04534,0.74649)(-4.03754,0.79976)%
(-4.02315,0.85123)(-4.00000,0.90000)(-3.98575,0.91654)(-3.96851,0.92540)(-3.94890,0.92807)%
(-3.92759,0.92607)(-3.90522,0.92089)(-3.88243,0.91404)(-3.85987,0.90702)(-3.83818,0.90134)%
(-3.81801,0.89850)(-3.80000,0.90000)(-3.77651,0.90964)(-3.75301,0.92537)(-3.72979,0.94493)%
(-3.70715,0.96607)(-3.68539,0.98654)(-3.66479,1.00408)(-3.64566,1.01646)(-3.62828,1.02140)%
(-3.61297,1.01667)(-3.60000,1.00000)(-3.57993,0.94705)(-3.57062,0.89344)(-3.56970,0.83967)%
(-3.57482,0.78623)(-3.58361,0.73361)(-3.59370,0.68230)(-3.60275,0.63279)(-3.60839,0.58557)%
(-3.60826,0.54115)(-3.60000,0.50000)(-3.59401,0.48713)(-3.58960,0.48909)(-3.58583,0.50307)%
(-3.58179,0.52623)(-3.57654,0.55576)(-3.56918,0.58883)(-3.55876,0.62261)(-3.54438,0.65428)%
(-3.52510,0.68102)(-3.50000,0.70000)(-3.46271,0.71615)(-3.42650,0.72720)(-3.39157,0.73371)%
(-3.35814,0.73626)(-3.32639,0.73541)(-3.29653,0.73173)(-3.26877,0.72578)(-3.24329,0.71813)%
(-3.22030,0.70935)(-3.20000,0.70000)%
%
\linethickness{0.008in}%%
\linethickness{0.016in}%%
\polyline(2.90000,0.80000)(2.89024,0.79091)(2.89056,0.77920)(2.89881,0.76494)(2.91288,0.74820)%
(2.93062,0.72906)(2.94992,0.70760)(2.96863,0.68389)(2.98464,0.65800)(2.99580,0.63001)%
(3.00000,0.60000)(2.99595,0.58147)(2.98515,0.56289)(2.96964,0.54415)(2.95146,0.52515)%
(2.93263,0.50579)(2.91519,0.48597)(2.90117,0.46558)(2.89261,0.44452)(2.89154,0.42270)%
(2.90000,0.40000)(2.93505,0.35154)(2.97552,0.31114)(3.01979,0.27733)(3.06624,0.24861)%
(3.11325,0.22350)(3.15920,0.20051)(3.20247,0.17815)(3.24144,0.15494)(3.27449,0.12939)%
(3.30000,0.10000)(3.30530,0.08692)(3.29858,0.08256)(3.28223,0.08474)(3.25867,0.09128)%
(3.23028,0.10000)(3.19947,0.10872)(3.16863,0.11526)(3.14018,0.11744)(3.11650,0.11308)%
(3.10000,0.10000)(3.08613,0.07481)(3.07664,0.04804)(3.07105,0.01987)(3.06886,-0.00949)%
(3.06959,-0.03986)(3.07275,-0.07104)(3.07786,-0.10284)(3.08443,-0.13506)(3.09197,-0.16751)%
(3.10000,-0.20000)(3.10512,-0.21375)(3.11276,-0.22574)(3.12237,-0.23631)(3.13343,-0.24582)%
(3.14539,-0.25461)(3.15772,-0.26303)(3.16988,-0.27143)(3.18134,-0.28016)(3.19156,-0.28957)%
(3.20000,-0.30000)(3.21197,-0.32149)(3.22127,-0.34471)(3.22873,-0.36887)(3.23520,-0.39322)%
(3.24151,-0.41698)(3.24850,-0.43938)(3.25701,-0.45965)(3.26787,-0.47703)(3.28192,-0.49073)%
(3.30000,-0.50000)(3.32217,-0.50515)(3.34436,-0.50609)(3.36640,-0.50314)(3.38807,-0.49661)%
(3.40920,-0.48682)(3.42960,-0.47408)(3.44907,-0.45871)(3.46742,-0.44103)(3.48446,-0.42136)%
(3.50000,-0.40000)(3.50950,-0.37600)(3.51266,-0.34544)(3.51108,-0.31110)(3.50633,-0.27579)%
(3.50000,-0.24230)(3.49367,-0.21343)(3.48892,-0.19196)(3.48734,-0.18071)(3.49050,-0.18245)%
(3.50000,-0.20000)(3.52218,-0.25107)(3.54045,-0.30563)(3.55610,-0.36313)(3.57040,-0.42299)%
(3.58465,-0.48465)(3.60012,-0.54753)(3.61811,-0.61108)(3.63990,-0.67472)(3.66677,-0.73788)%
(3.70000,-0.80000)(3.71020,-0.81107)(3.72026,-0.81152)(3.73023,-0.80357)(3.74013,-0.78946)%
(3.75000,-0.77141)(3.75987,-0.75164)(3.76977,-0.73239)(3.77974,-0.71588)(3.78980,-0.70434)%
(3.80000,-0.70000)(3.82151,-0.70575)(3.84595,-0.72088)(3.87208,-0.74223)(3.89869,-0.76664)%
(3.92456,-0.79093)(3.94846,-0.81195)(3.96918,-0.82654)(3.98549,-0.83152)(3.99617,-0.82372)%
(4.00000,-0.80000)(3.99315,-0.73185)(3.97475,-0.66634)(3.94803,-0.60318)(3.91624,-0.54204)%
(3.88260,-0.48260)(3.85036,-0.42456)(3.82274,-0.36759)(3.80299,-0.31139)(3.79433,-0.25563)%
(3.80000,-0.20000)(3.81301,-0.17148)(3.83120,-0.16507)(3.85404,-0.17650)(3.88101,-0.20155)%
(3.91158,-0.23597)(3.94522,-0.27551)(3.98141,-0.31592)(4.01962,-0.35297)(4.05933,-0.38241)%
(4.10000,-0.40000)(4.11873,-0.40251)(4.13557,-0.40061)(4.15051,-0.39492)(4.16352,-0.38604)%
(4.17458,-0.37458)(4.18368,-0.36116)(4.19079,-0.34638)(4.19589,-0.33085)(4.19897,-0.31519)%
(4.20000,-0.30000)(4.19720,-0.28118)(4.18960,-0.26284)(4.17840,-0.24473)(4.16480,-0.22653)%
(4.15000,-0.20799)(4.13520,-0.18880)(4.12160,-0.16869)(4.11040,-0.14738)(4.10280,-0.12457)%
(4.10000,-0.10000)(4.10376,-0.07899)(4.11382,-0.05866)(4.12834,-0.03882)(4.14547,-0.01933)%
(4.16338,0.00000)(4.18021,0.01933)(4.19413,0.03882)(4.20330,0.05866)(4.20587,0.07899)%
(4.20000,0.10000)(4.18231,0.13060)(4.16118,0.15890)(4.13687,0.18487)(4.10964,0.20850)%
(4.07977,0.22977)(4.04751,0.24865)(4.01314,0.26513)(3.97692,0.27920)(3.93912,0.29083)%
(3.90000,0.30000)(3.87544,0.29892)(3.85009,0.28807)(3.82463,0.27036)(3.79973,0.24871)%
(3.77606,0.22606)(3.75431,0.20533)(3.73516,0.18943)(3.71927,0.18129)(3.70732,0.18384)%
(3.70000,0.20000)(3.69421,0.25975)(3.69895,0.31966)(3.71158,0.37970)(3.72947,0.43983)%
(3.75000,0.50000)(3.77053,0.56017)(3.78842,0.62030)(3.80105,0.68034)(3.80579,0.74025)%
(3.80000,0.80000)(3.79218,0.81625)(3.77920,0.81796)(3.76197,0.80838)(3.74141,0.79079)%
(3.71846,0.76846)(3.69403,0.74465)(3.66904,0.72263)(3.64443,0.70567)(3.62111,0.69704)%
(3.60000,0.70000)(3.56103,0.72803)(3.53034,0.76610)(3.50641,0.81061)(3.48772,0.85795)%
(3.47273,0.90453)(3.45993,0.94675)(3.44778,0.98101)(3.43476,1.00370)(3.41934,1.01123)%
(3.40000,1.00000)(3.37756,0.96497)(3.37009,0.92307)(3.37382,0.87564)(3.38504,0.82400)%
(3.40000,0.76948)(3.41496,0.71339)(3.42618,0.65708)(3.42991,0.60186)(3.42244,0.54906)%
(3.40000,0.50000)(3.38626,0.48682)(3.37252,0.48741)(3.35826,0.49938)(3.34295,0.52032)%
(3.32608,0.54783)(3.30712,0.57952)(3.28556,0.61298)(3.26086,0.64581)(3.23252,0.67562)%
(3.20000,0.70000)(3.15926,0.72386)(3.12067,0.74516)(3.08431,0.76371)(3.05028,0.77932)%
(3.01867,0.79178)(2.98956,0.80090)(2.96305,0.80647)(2.93922,0.80832)(2.91818,0.80622)%
(2.90000,0.80000)%
%
\linethickness{0.008in}%%
\linethickness{0.04in}%%
\polyline(-4.00000,0.90000)(-4.00000,2.10000)%
%
\linethickness{0.008in}%%
\linethickness{0.04in}%%
\polyline(3.80000,0.70000)(4.00000,1.90000)%
%
\linethickness{0.008in}%%
\end{picture}}%}}
\end{layer}\par
\vspace{11\baselineskip}

下のQRコードからリンクへ飛んでみて。\\

%% upLaTeXでは，図の挿入にBBdataが必須
\includegraphics[bb=10 10 220 220,width=50mm,clip]{kamirinQR_157881.png}\\

\url{https://kubo-yuge.github.io/kamirin22/kamirinOspSample201003json.html}\\

\newpage

\end{multicols}

\end{document}
